\documentclass{article}
\usepackage{hyperref}
\hypersetup{
    urlcolor=blue,
    citecolor=blue,
    colorlinks=true
%    colorlinks=true
}

%    filecolor=magenta,      
%    urlcolor=cyan,
%    pdftitle={Overleaf Example},
%    pdfpagemode=FullScreen,
%    }
\usepackage{comment}
\usepackage{tcolorbox}
\usepackage{capt-of}
\usepackage{amsmath}
\usepackage{algorithm}
\usepackage{algpseudocode}
\usepackage{cancel}
\usepackage{amsfonts} 
\usepackage{textcomp}
\usepackage{amssymb}
\usepackage{subfig}
\usepackage{caption}
\usepackage{graphicx}
\usepackage{wrapfig}
\usepackage{float}
\usepackage{cancel}
\usepackage{booktabs}
\usepackage{lipsum}
\usepackage{enumerate}
\usepackage{authblk}
\usepackage{makecell}
\usepackage[a4paper, total={6.5in, 9.5in}]{geometry}
\usepackage[utf8]{inputenc}
\usepackage[]{apacite}
\bibliographystyle{apacite}
\usepackage{lineno}
\usepackage{xcolor}
\linenumbers
\DeclareMathOperator*{\argminA}{arg\,min} 
\title{Semantic representations in episodic memory enhance recall and compositional consolidation}
\author[1]{Albert Albesa-González}
\author[1]{Claudia Clopath} 
\affil[1]{Department of Bioengineering, Imperial College London, London, UK}
\begin{document}
\maketitle
\begin{abstract}
\noindent Episodic and semantic memory are classically thought to play distinct roles: episodic memory encodes unique experiences, while semantic memory generalizes across them. Current conceptualizations of episodic and semantic memory interactions emphasize a one-way consolidation from episodic traces in the medial temporal lobe (MTL) to semantic knowledge in neocortex (CTX). However, this tradition has left largely unexplored how semantic memory may affect episodic encoding of new memories. Here we introduce a cognitive model in which the code used in episodic memories shifts from purely sensorial to including explicit semantic representations of stored events. Simultaneously, we propose a computational circuit model of how such a cognitive strategy could be implemented in the brain using biologically-plausible learning rules. We show that increased sparsity during replay enables neocortex to extract compositional structure from overlapping episodes, which creates a dictionary of inter-connected concepts in semantic memory. Furthermore, we show that spontaneous activity in neocortical areas can imprint the abstracted representations into the medial temporal lobe, giving rise to concept-like cells. This bidirectional interaction improves episodic recall and replay fidelity, and  facilitates the consolidation of higher-order representations based on previous semantic knowledge. The model accounts for behavioural advantages of schema-congruent learning, the emergence of concept neurons, and enhanced memory performance for semantically familiar stimuli. Together, our results provide a mechanistic account of how episodes and concepts reinforce each other, extending standard consolidation theories and suggesting a cooperative framework where semantic knowledge scaffolds episodic encoding, which in turn favours compositional abstraction.
\end{abstract}
\section*{Introduction}
The distinction between episodic and semantic memory has been central to cognitive neuroscience since its first formalization \shortcite{tulving_1972}. In its original view, episodic memory captures individual experiences, while semantic memory accumulates generalized knowledge that emerges from repeated exposure to similar episodes \shortcite{squire_2004}. One influential computational explanation for this division is provided by the Complementary Learning Systems (CLS) framework. This theory posits that a fast-learning system (the hippocampus and the rest of the medial temporal lobe, MTL) specializes in rapidly encoding episodic memories. In contrast, a slow-learning neocortical (CTX) system works in concert with the MTL to consolidate and generalize memories \shortcite{mcclelland_1995, oreilly_2014}.
\newline\newline
Complementary Learning Systems and other consolidation models have historically emphasized a unidirectional transfer of knowledge from episodic traces (MTL) to semantic schemas (CTX) \shortcite{moscovitch_2021}. In this sense, models of the opposite (cortex-to-hippocampus teaching) have remained widely unexplored, as recently pointed out in \shortciteA{liu_2024}. However, humans often remember well-structured material more reliably than random or unfamiliar content, suggesting that top-down, schema-like influences can guide the encoding and retrieval of new events \shortcite{bartlett_1932, alba_1983, lin_2021}. On the neurobiological side, research has identified so-called \textit{concept cells} in the human MTL. These are neurons that fire selectively for particular semantic categories or even specific concepts (such as the famously named \textit{Jennifer Aniston neuron}) \shortcite{quiroga_2005, quiroga_2012, courellis_2024}. Furthermore, regions in the medial temporal lobe such as the hippocampus and the entorhinal cortex have been shown to contain neurons coding for latent elements in the environment. These especially include latents related to space \shortcite{o_1971, fyhn_2004, moser_2008}, but also other forms of complex task structure \shortcite{aronov_2017, sun_2020, zhao_2022}. Finally, the existence of long-range prefrontal-to-hippocampal inhibitory connections suggests that abstract information represented in neocortical areas can be directly conveyed to subcortical regions \shortcite{malik_2022}, with these regions exerting top-down control of subcortical areas during sleep \shortcite{swanson_2020, shin_2024}, typically associated with consolidation. All in all, these results suggest that abstract semantic information can be strongly represented in the Medial Temporal Lobe regions, despite its canonical role in episodic memory.
\newline\newline
While views challenging a strict dichotomy between semantic and episodic memory are almost as old as the semantic-episodic distinction itself \shortcite{mckoon_1986, greenberg_2010, gentry_2024}, we still lack a mechanistic understanding of how or why these two memory systems interact. First, most consolidation models assume that replay from hippocampus to cortex is a faithful reinstatement of cortical activity, often overlooking the biological reality that such replays are inherently noisy and the resulting cortical representations are actively constructed rather than simply copied \shortcite{swanson_2020}. Furthermore, highly semantized pre-existing representations are usually assumed \shortcite{saxe_2019, chrysanthidis_2022} sidestepping how such mappings might arise through local, biologically plausible learning rules in the first place. Another key challenge is reconciling the presence of semantic representations in hippocampus and rest of MTL with attractor models of episodic memory. Semantic representations naturally induce overlapping patterns, which have been known for a very long time to reduce memory capacity. For this reason, consolidation models many times assume hash-like codes or other forms of perfect pattern separation, even when the episodes are conceptually highly overlapping \shortcite{sun_2023, chandra_2025}.
\newline\newline
Here, we introduce a model of episodic and semantic memory interaction with two key processes: (i) episodic-to-semantic consolidation and (ii) semantic-to-episodic consolidation. Episodic-to-semantic consolidation is associated with the standard gradual transfer of information of episodic regularities into semantic memory. Conversely, episodic-to-semantic consolidation transfers these regularities into the episodic code, actively shaping the representations of future episodes based on previously formed abstractions. This cognitive model is presented together with a biologically-motivated circuit model of how such a cognitive algorithm could be implemented. The circuit, based on canonical models of subcortical-cortical consolidation, learns using only local signals via Hebbian and homeostatic plasticity. We use compositional input with a pre-defined statistical structure, which the neocortical component of the model is able to extract and decompose. These compositional representations are then transferred into the episodic memory system (Medial Temporal Lobe), increasing the semantic content of episodic representations. Our model can explain the formation of concept-like cells in subcortical areas \shortcite{quiroga_2005, quiroga_2012}, the advantage of blocked learning in humans \shortcite{flesch_2018}, and the better recall of memories with semantic content \shortcite{lin_2021}. Altogether, our model suggests that semantic memory can enhance episodic memory by transferring into it semantic representations. Furthermore, the presence of these semantic representations in episodic encoding facilitate the consolidation of concepts of higher complexity in semantic memory, supporting learning of higher-order abstractions. These results extend previous models of unidirectional transfer of information from hippocampus to cortex, pointing to a cooperative mechanism where both memory systems enhance each other.  
\newpage
\section*{Results}
\subsection*{A cognitive-circuit model for semantic structure in episodic memory}
We introduce a cognitive model of semantic and episodic memory, together with a circuit implementation (Fig. \ref{fig:diagram}). The model starts from a simple psychological observation: episodic memories often contain both sensory detail and semantic content. For example, let's consider multiple episodes of someone having their favourite food: \textit{pizza}. These may differ in appearance, smell, taste, and context, but still share semantic structure such as \textit{pizza}, \textit{cheese}, and \textit{tomato}. Our cognitive model captures this by positing that episodic traces contain two components (Fig. \ref{fig:diagram}a): a sensory component, preserving event-specific perceptual detail, and a semantic component, representing reusable concepts shared across events.
\newline\newline
At the circuit level, however, this raises a computational problem. The medial temporal lobe (MTL), a region crucial for episodic memory, has been shown to contain concept-like information. Yet classical attractor models achieve maximal capacity when stored patterns are decorrelated or approximately orthogonal, whereas semantic representations introduce structured overlap between episodes. How can an episodic memory system tolerate this overlap without collapsing into interference? To address this question, we build a complementary circuit model to study how semantic information can coexist with episodic storage in MTL. Following standard mappings, neocortex (CTX) stores semantic memory and MTL stores episodic memory (Fig. \ref{fig:diagram}c). MTL is divided into MTL-sensory, driven by sensory input, and MTL-semantic, driven by higher-order cortical input. These populations are recurrently connected, allowing sensory and semantic representations to be bound into individual episodic traces.
\newline\newline
A second issue concerns how cortical representations are learned from replay. A common framing is that cortical reactivations are used as a teaching signal, which can then be back-propagated. However, this often leaves implicit how the reactivated episodic pattern interacts with cortical activity. Which cortical neurons provide the input? Which ones define the target? Which ones remain plastic? Even biologically motivated alternatives to backpropagation, such as Contrastive Hebbian Learning or Predictive Coding, still require a structured distinction between input, hidden, and output components during learning.
\newline\newline
We avoid this by proposing replay as a mechanism for dictionary learning (Technical Box \ref{box:replay}). In this view, the role of replay is to provide episodic patterns that CTX must interpret.
During \textit{episodic replay}, spontaneous MTL activity evolves through recurrent connections that contain statistics accumulated across episodes. The replayed pattern is not a verbatim copy of awake experience, but a generalized episodic pattern. When projected to CTX, this pattern either refines existing representations of similar patterns or recruits new neurons if no previous representation is sufficiently close. This is achieved through local Hebbian plasticity and competition, avoiding both backpropagation-like credit assignment and the need to predefine cortical targets.
\newline\newline
Once CTX has learned this vocabulary, the same logic can operate in reverse. If recurrent MTL dynamics expose episodic structure to CTX, recurrent CTX dynamics can expose semantic structure to MTL. We call this process \textit{semantic replay} because, like episodic replay, is internally generated, begins from random activity, and recovers attractor states through recurrent dynamics. However, because these attractors are shaped by CTX recurrent connections, they reflect semantic rather than episodic structure. Projecting them to MTL-semantic provides a local learning signal through which concept-like representations can enter the episodic circuit.
\newline\newline
Together, episodic and semantic replay implement the two transformations proposed by the cognitive model (Fig. \ref{fig:diagram}b,d). Episodic replay supports \textit{episodic-to-semantic consolidation}, allowing CTX to extract a cortical dictionary from episodic structure. Semantic replay supports \textit{semantic-to-episodic consolidation}, allowing MTL-semantic to incorporate the learned cortical vocabulary. We next validate this learning algorithm in simulations of the circuit model. We then use the model to ask how can semantic overlap interact in a beneficial way with episodic attractor dynamics, supporting memory and further consolidation.
\begin{figure}[h!]
    \centering
    \includegraphics[width=\linewidth]{Figures/Figure_1.png}
    \caption{\textbf{We propose a cognitive model in which semantic memory shapes representations in episodic memory, and an associated circuit implementation.} \textbf{a}: Cognitive model of episodic (red) and semantic (green) memory. Episodic memories integrate both sensory information (orange) and semantic content (green), while semantic memory extracts regularities across episodes. 
\textbf{b}: Schematic of bidirectional consolidation. Episodic-to-semantic consolidation transfers regularities from episodic memory to semantic memory, while semantic-to-episodic consolidation builds semantic representations in episodic memory.  
\textbf{c}: Circuit implementation of the model. Neocortex (CTX) stores semantic memories and Medial Temporal Lobe (MTL) stores episodic memories. MTL is split into neurons primarily receiving input from sensory cortex (MTL-sensory, coding for low-level sensory features) and neurons primarily receiving input from neocortex (MTL-semantic, coding for high-level abstract features). 
\textbf{d}: Mapping of cognitive consolidation processes onto replay mechanisms. Episodic replay corresponds to CTX$\leftarrow$MTL interactions (MTL-to-CTX connections trained during replay), while semantic replay corresponds to MTL-semantic$\leftarrow$CTX interactions (semantic assemblies in CTX driving concept representations in MTL-semantic).}
    \label{fig:diagram}
\end{figure}
\begin{comment}
\begin{tcolorbox}[colback=blue!5!white, colframe=blue!75!black, title=Technical Box - Input Generation]
\textbf{Intuition:} We generate compositional \emph{episodes} by combining one concept from each of two attributes. The different letters denote attributes (categories) and subscripts denote instances within each attribute. For example, $A_i$ could represent different visual cues and $B_j$ different auditory cues. An \emph{episode} is the pair $\{A_i,B_j\}$, and episodes are sampled from a joint distribution $p(A_i,B_j)$. Importantly, this joint distribution induces an asymmetric \emph{semantic structure} between concepts via conditional probabilities:
\[
\mathcal{SS}_{xy} \;=\; p(x \in e \mid y \in e),
\]
which we call ``semantic'' because it defines the meaning of a concept by the concepts that reliably co-occur with it. For instance, the semantics of \textit{summer} are given by concepts that are likely given \textit{summer} (sun, beach) and unlikely given \textit{summer} (snow, cold). 
\newline\newline
\textbf{Relevant details:}
\begin{itemize}
    \item \textbf{Sampling episodes within a day.} A day of sensory input has length $T_\textrm{day}$ timesteps and is partitioned into $T_\textrm{day}/T_\textrm{episode}$ blocks. At the start of each block, we sample one episode
    \[
    e=\{A_i,B_j\} \sim p(A_i,B_j),
    \]
    and hold its underlying concepts fixed for $T_\textrm{episode}$ timesteps.
    
    \item \textbf{Mapping an episode to a sensory activity pattern.} Each concept $c\in A\cup B$ is assigned a dedicated subpopulation of sensory neurons, $\mathrm{SEN}_c$ (Fig.~XXX). Given a sampled episode $e=\{A_i,B_j\}$, we generate an initial binary sensory vector $X^{\mathrm{SEN}}$ by activating exactly the neurons belonging to the two concept-specific subpopulations:
    \[
    X^{\mathrm{SEN}}_n =
    \begin{cases}
    1 & \text{if } n\in \mathrm{SEN}_{A_i}\ \text{or}\ n\in \mathrm{SEN}_{B_j},\\
    0 & \text{otherwise.}
    \end{cases}
    \]
    Operationally, $\mathrm{SEN}_{A_i}$ and $\mathrm{SEN}_{B_j}$ correspond to two concatenated ``blocks'' (one for attribute $A$, one for attribute $B$), with each block containing a fixed number of neurons per concept (Fig.~XXXX). This yields a structured, sparse sensory pattern whose active set explicitly encodes the current episode concepts.
    
    \item \textbf{Within-episode variability (swap noise).} To model variability across repeated presentations of the same episode, we inject controlled stochasticity at every timestep. Starting from the base pattern above, we uniformly sample $N_\textrm{swaps}$ active neurons and $N_\textrm{swaps}$ inactive neurons and flip their states (Fig.~XXXXX). This produces a noisy realization while preserving overall activity level (sparsity is unchanged because the number of $1\!\to\!0$ flips equals the number of $0\!\to\!1$ flips).
    
    \item \textbf{Temporal sequence.} Repeating the above over blocks produces a day-long tensor $X_\textrm{input}\in\{0,1\}^{(N_\textrm{days},T_\textrm{day}, N_\textrm{SEN}})$ in which (i) the episode identity changes every $T_\textrm{episode}$ timesteps and (ii) noise is independently resampled at each timestep.
\end{itemize}
\end{tcolorbox}
\begin{tcolorbox}[colback=blue!5!white, colframe=blue!75!black, title=Technical Box - Network Dynamics and Local Learning]
\textbf{Intuition:} All regions are composed of binary neurons operating under strong competition. At any given timestep, only a fixed number of neurons are allowed to be active. This \emph{K-winners-take-all} (K-WTA) mechanism enforces sparse, distributed representations and induces competition between neurons to represent incoming patterns. Learning is purely local and Hebbian, but combined with divisive normalization (homeostasis). The interaction between competition and normalization makes synaptic weights converge to scaled conditional probabilities, so that connectivity reflects statistical structure in the environment.
\newline
\textbf{Details: } Each region receives a pre-activation synaptic input:
\begin{equation}
    \hat{X}^\textrm{region} = W^{\textrm{region}\leftarrow\textrm{pre}}\cdot\hat{X}^\textrm{region}
\end{equation}
which is then mapped into binary activations via a top-K operation (therefore implementing \textit{K-winners-take-all} dynamics):
\begin{equation}
    X^\textrm{region} = \textrm{top-}K\big(\hat{X}^\textrm{region}\big)
\end{equation}
The basic form of learning is Hebbian: 
\begin{equation}
    \Delta W_{ij}^{\textrm{post}\leftarrow\textrm{pre}} = \lambda_{ij}X^\textrm{post}_iX^\textrm{pre}_j
\end{equation}
which is complemented with divisive normalization. When synaptic weights are divided by the sum of incoming weights $\sum_iW_{ij}$, we call this process \textit{incoming homeostasis}. Under this form of learning, the weights behave as an exponential moving average of conditional probabilities, converging to:
   \begin{equation}
   W_{ij}^{\textrm{pre}\leftarrow\textrm{post}, \; \textrm{incoming}}(\infty) \equiv \frac{W_\textrm{max-in}^{\textrm{pre}\leftarrow\textrm{post}}}{K^\textrm{pre}}p(X^\textrm{pre}_j = 1| X^\textrm{post}_i = 1)
 \end{equation}
and when these are divided by the sum of outgoing weights $\sum_iW_{ij}$, the fixed points are:
  \begin{equation}
     W_{ij}^{\textrm{pre}\leftarrow\textrm{post}, \; \textrm{outgoing}}(\infty) \equiv \frac{W_\textrm{max-out}^{\textrm{pre}\leftarrow\textrm{post}}}{K^\textrm{post}}p(X^\textrm{post}_i = 1| X^\textrm{pre}_j = 1)
 \end{equation}
 \textbf{Motivation:} The basic neuronal and learning dynamics are inspired by classic models of competitive learning and sparse coding CITE. The implementation details are heavily based on recent work 
\end{tcolorbox}
\end{comment}
\newtcolorbox[auto counter]{technicalbox}[2][]{
  colback=blue!5!white,
  colframe=blue!75!black,
  title=Technical Box~\thetcbcounter: #2,
  label=#1
}
\clearpage
\begin{technicalbox}[box:replay]{Replay for Dictionary Learning}
Instead of viewing replay as perfect reinstatement of stored episodes, we propose it can be seen as a mechanism for dictionary learning, where a \textit{post} region learns relevant \textit{atoms} describing a \textit{pre} region based on presynaptic recurrent connections. The mechanism is as follows:
\newline\newline
\textbf{Generate a presynaptic pattern} from noise via attractor dynamics:
\[
X^\mathrm{pre}_{0} \sim \mathcal{N}(0,I_{N_\mathrm{pre}})
\quad \longmapsto \quad
X^\mathrm{pre}_{n+1}
=
\operatorname{top\text{-}K_R^{\mathrm{replay}}}
\left(
W^{\mathrm{pre}\leftarrow \mathrm{pre}} X^\mathrm{pre}_{n}
\right).
\]
\textbf{Project to postsynaptic layer}
\[
X^\mathrm{post}
=
\operatorname{top\text{-}K^\textrm{post}}
\left(
W^{\mathrm{post}\leftarrow\mathrm{pre}}X^\mathrm{pre}
+
b_\mathrm{IM}^\textrm{post}\mathrm{IM^\textrm{post}}
\right),
\]
\textbf{Update dictionary} For postsynaptic active neurons $i$ ($X^\mathrm{post}_i$ = 1).
\begin{itemize}
    \item If neuron $i$ is immature ($\textrm{IM}^\textrm{post}_i$ = 1): learn new atom
    \begin{itemize}
        \item Create one shot receptive field $W_i^{\textrm{post}\leftarrow\textrm{pre}} = X^\textrm{pre}$
        \item Make neuron mature ($\textrm{IM}^\textrm{post}_i$ = 1)
    \end{itemize}
    \item If neuron $i$ is mature  ($\textrm{IM}^\textrm{post}_i$ = 0): update existing atom
    \begin{itemize}
        \item Receptive field is pushed in the direction of $X^\textrm{pre}$
        \[W_{ij}^{\textrm{post}\leftarrow\textrm{pre}} \longmapsto
p(X^\mathrm{pre}_j=1 \mid X^\mathrm{post}_i=1)\]
    \end{itemize}
\end{itemize}
\end{technicalbox}
\begin{comment}
\begin{tcolorbox}[colback=blue!5!white, colframe=blue!75!black, title=Technical Box - General Learning Algorithm]
\textbf{Intuition:} 
Learning alternates between two regimes: \emph{wake} (input-driven encoding) and \emph{sleep} (replay-driven abstraction). During wake, sensory input drives the network and updates recurrent synapses locally: MTL rapidly stores episodic traces, while CTX slowly accumulates statistical regularities. During sleep, internally generated replay transfers structure across regions. Depending on the training phase, replay progressively shifts from purely sensory episodic content to increasingly semantic and compositional structure.
\newline\newline
\textbf{Wake (online encoding):}
For each day, sensory input is processed sequentially:
\begin{enumerate}
    \item Sensory activity drives MTL-sensory and (initially random) MTL-semantic populations.
    \item MTL activity projects to CTX, where sparse competition selects a cortical representation.
    \item Recurrent synapses are updated locally:
    \begin{itemize}
        \item MTL recurrent connections store the current episode (rapid episodic encoding).
        \item CTX recurrent connections are updated with a slow learning rate, gradually capturing co-occurrence statistics (semantic structure).
    \end{itemize}
\end{enumerate}

At early stages, MTL-semantic activity is not yet driven by CTX and episodic traces contain only sensory content. After cortical structure has formed (Phase A), CTX begins to drive MTL-semantic, enriching episodic representations with semantic components. In later stages (Phase B), episodic traces contain both sensory and semantic structure, allowing higher-order regularities to be stored.
\newline\newline
\textbf{Sleep (replay-driven reorganization):}
Sleep consists of two replay loops:

\begin{itemize}
    \item \emph{Episodic replay (MTL $\rightarrow$ CTX):} 
    Recurrent MTL dynamics generate structured patterns which train CTX feed-forward synapses. Early in training, replay involves only sensory MTL content. After Phase B, replay draws from full MTL representations (sensory + semantic), and the semantic load can vary, enabling compositional abstraction.

    \item \emph{Semantic replay (CTX $\rightarrow$ MTL-semantic):} 
    Once cortical recurrent structure is established (after Phase A), CTX replay trains MTL-semantic representations, creating pointers to cortical statistical regularities.
\end{itemize}

Each replay step follows the same local rule: recurrent pattern completion in the \emph{pre} region, projection to the \emph{post} region, and Hebbian learning with homeostasis. The difference across phases lies only in which region replays to which.
\newline\newline
\textbf{Phase structure:}
\begin{itemize}
    \item \textbf{Phase A:} MTL stores sensory episodes during wake. Sleep replay trains CTX to extract sensory regularities.
    \item \textbf{Phase B:} CTX representations begin to drive MTL-semantic during wake. Sleep additionally transfers cortical structure back to MTL-semantic.
    \item \textbf{Late phase:} Episodic traces contain both sensory and semantic components. Replay operates over this richer space, enabling higher-order abstraction.
\end{itemize}
\newline
\textbf{Motivation:} There are some unusual choices in the general learning algorithm. For instance, the cortical generative model (embedded in cortico-cortical synapses) is not trained during sleep. Instead, sleep re-organizes cortical and medial temporal representations. A similar approach has been proposed in the wake-sleep library learning algorithm for program synthesis of DreamCoder CITE. Furthermore, \textit{episodic replay} can be mapped to the \textit{Abstraction} phase of DreamCoder, with both models creating new pointers that represent regularities present across episodes. In a similar fashion, \textit{semantic replay} can be compared to the corresponding \textit{Dreaming} (program recognition) sleep phase, with our network learning how to represent newly encoded episodes using the pointers discovered during episodic replay. Add summary regarding semantic replay biological plausibility. 
\end{tcolorbox}
\end{comment}
\newpage
\subsection*{Episodic replay can implement dictionary learning in semantic memory}
We first consider episodic-to-semantic consolidation (Fig. 2a) as a form of dictionary learning carried out by neocortex. In dictionary learning, complex inputs are represented by combining a small number of reusable elements (typically called \textit{atoms}, Technical Box \ref{box:replay}). Here, we refer to these learned elements as \textit{cortical words}: reusable cortical representations that capture recurring parts of episodic experience. In this view, episodic replay in MTL does not simply reinstate full episodes. Rather, sparse replay proposes recurring subpatterns embedded across many episodes. CTX then either matches each replayed candidate to a previously learned cortical word, allowing the corresponding receptive field to be refined, or treats it as sufficiently novel to recruit immature cortical neurons and establish a new receptive field. In this way, repeated episodic replay gradually builds a cortical dictionary of reusable semantic components.
\newline\newline
To test this idea, we generated episodes by combining pairs of latent factors $(A_i,B_j)$, with $A_i \in \{A_1,\dots,A_5\}$ and $B_j \in \{B_1,\dots,B_5\}$ (Fig. \ref{fig:phase_a}b; see Methods). Each latent factor (e.g. $A_2$) corresponds to a semantically coherent sensory feature, such as the tomato or cheese element of a pizza. The distribution over pairings was structured: matched pairs, such as $(A_3,B_3)$, accounted for 50\% of episodes, whereas all mismatched pairings jointly accounted for the remaining 50\%. Thus, although every $A_i$ could in principle appear with every $B_j$, factors with matching indices co-occurred more often. For example, $B_1$ was statistically more associated with $A_1$ than with $A_2$.
\newline\newline
During wake, episodic inputs are stored in MTL through Hebbian plasticity. During sleep, sparse pattern completion (Technical Box X) in MTL recovers subpatterns shared across episodes, rather than reinstating complete episodic states (Fig. \ref{fig:phase_a}c). These replayed subpatterns are projected to CTX through initially random CTX$\leftarrow$MTL connections (Fig. \ref{fig:phase_a}d, $t=0$). If a replayed pattern is sufficiently close to an existing cortical receptive field, the corresponding cortical word is reactivated and refined. Otherwise, because immature cortical neurons have higher excitability (Technical Box X), the pattern recruits a previously unused population and forms a one-shot cortical representation (Fig. \ref{fig:phase_a}d, $t=200$). Subsequent replay of related patterns then drives slow Hebbian learning and homeostatic competition, progressively refining these receptive fields so that cortical neurons become tuned to stable, de-noised components of episodic experience (Fig. \ref{fig:phase_a}d,e).
\newline\newline
We confirmed this process by measuring the mutual information between cortical activity and the underlying \textit{simple concepts}, $A_i$ and $B_j$. Cortical activity rapidly became maximally informative about these latent factors, reaching a normalized mutual information of 1 (Fig. \ref{fig:phase_a}f). \textit{Complex concepts} of the form $A_iB_j$ are considered in later sections. We also found that the probability of learning a cortical word increased monotonically with the probability that its corresponding simple concept appeared in an episode (Fig. \ref{fig:phase_a}g). Thus, this mechanism can be interpreted as a reward-independent form of prioritized replay: concepts that are more likely to describe future episodes are also more likely to be consolidated as cortical words.
\newline\newline
Recurrent CTX$\leftarrow$CTX connections support a complementary form of dictionary learning. When multiple cortical neurons become tuned to the same replayed candidate, they also become strongly recurrently connected, forming a cortical assembly for that word. This corresponds to the standard replay-mediated formation of recurrent cortical assemblies, in which re-activated neurons are bound into a stable attractor-like representation. In parallel, slow plasticity updates recurrent connections across the broader cortical population. Under Hebbian plasticity with outgoing homeostasis, these recurrent weights converge to the conditional firing structure between neurons \cite{albesagonzalez_2025}. Because cortical neurons are tuned to individual concepts, recurrent CTX connectivity captures how likely one concept is to co-occur with another, thereby encoding the associative structure of the dictionary itself (Fig. \ref{fig:phase_a}h). Thus, while feedforward CTX$\leftarrow$MTL weights learn the dictionary entries---what each cortical word encodes--- recurrent CTX$\leftarrow$CTX connections learn the relations between entries---how each word is situated within the broader semantic structure.
\newline\newline
Together, these results suggest that semantic memory can emerge from episodic replay through cortical dictionary learning. Replay extracts candidate words from overlapping episodes; CTX stabilizes them as reusable semantic units; and recurrent cortical connectivity organizes them according to their statistical relations. In the pizza example, cortical words encode ingredients such as cheese or tomato, whereas corticocortical connections encode how strongly one ingredient predicts another.
\begin{figure}
    \centering
    \includegraphics[width=\linewidth]{Figures/Figure_2.png}
\caption{\textbf{Sparse episodic replay creates compositional representations in semantic memory.} \textbf{a}: Episodic-to-semantic consolidation via episodic replay. Cognitive model (left) and circuit implementation (middle/right). Sensory episodes (no semantics extracted yet) are stored in MTL-sensory, replayed sparsely during sleep to uncover overlapping sub-patterns, and then projected to CTX, where compositional semantic structure is extracted in CTX$\leftarrow$MTL (replay) and CTX$\leftarrow$CTX synapses (awake).  
\textbf{b}: Example neuronal activity during wake. Input episodes are sampled following a protocol as described in \shortciteA{albesagonzalez_2025}. To determine MTL-sensory activity, two different concepts (one $A_i \in A$ and one $B_j \in B$) are simultaneously sampled. 50\% of episodes are of form $(A_i, B_{j=i})$, and 50\% $(A_i, B_{j\neq i})$. This makes $A_i$ more semantically related to $B_i$ than the rest of $B$'s. 10 neurons in MTL-sensory code for each concept, and each episode lasts 5 timesteps. At each timestep, $N_\textrm{swaps}= 4$ random neurons that code for the present concepts are turned off, and  $N_\textrm{swaps}= 4$ random neurons that code for non-present concepts are turned on (this introduces noise with a consistent signal-to-noise ratio). CTX activity is driven by random initial CTX$\leftarrow$MTL-sensory connections.
\textbf{c}: Example neuronal activity during episodic replay. Sparse replay in MTL-sensory reveals concept-specific sub-patterns. These create structure in CTX$\leftarrow$MTL-sensory connections via one-shot followed by slow learning (Fig. \ref{fig:phase_a}d).}
    \label{fig:phase_a}
\end{figure}
\begin{figure}[t]
  \ContinuedFloat
  \captionsetup{list=off} % don't add a new list-of-figures entry
  \caption{\textit{(continued)} \textbf{d}: Example CTX neuron receptive field (CTX$\leftarrow$MTL synapses) over training. Images show weights associated to each presynaptic MTL-sensory neuron rearranged in a rectangular pattern according to the concept they code for (each row represents weights from neurons coding for a particular concept in MTL-sensory). Initially random connectivity ($t = 0$) is one-shot driven to represent a replayed presynaptic pattern ($t=200$). Then, slow Hebbian and homeostatic plasticity drive statistical learning, uncovering receptive fields that represent semantic components ($t = 200-90000$). Right: weight dynamics for all synapses of the same neuron. \textbf{e}: Final CTX$\leftarrow$MTL weight matrix. Neurons in CTX become tuned to individual concepts ($A_1$–$A_5$, $B_1$–$B_5$), forming a dictionary of abstract components.
\textbf{f}: Example activity during late wake. CTX neurons now respond selectively to semantic components, showing compositional coding aligned with MTL-sensory input.  
\textbf{g}: Final CTX$\leftarrow$CTX weight matrix (left) and relationship between synaptic strength and semantic structure (right). Recurrent connectivity captures conditional firing probabilities across concepts, organizing into block-like structure encoding semantic associations.}
\end{figure}
\newpage
\subsection*{Semantic replay can transfer the learned vocabulary into episodic memory}

We next asked whether spontaneous cortical activity can transfer the learned cortical dictionary back into episodic memory (semantic-to-episodic consolidation; Fig. \ref{fig:phase_b}a, left). We refer to this activity as \textit{semantic replay} because it parallels episodic replay: it is not directly driven by sensory input, and activity patterns are recovered through attractor dynamics from random initial states. However, here the replayed content comes from recurrent connections representing \textit{semantic} information, as opposed to episodics. In the pizza example, once ingredients such as tomato and cheese have been consolidated as cortical words, semantic replay can imprint these words onto MTL-semantic. Thus, before semantic-to-episodic consolidation, an episode containing tomato sauce may be encoded only through its sensory features; after consolidation, MTL-semantic contains neurons that explicitly signal the presence of tomato, analogous to MTL concept cells. At the circuit level, this corresponds to MTL-semantic neurons developing receptive fields that are selective for cortical assemblies representing individual concepts (Fig. \ref{fig:phase_b}a, right).
\newline\newline
To support this process, we now equip the simulated model with semantic replay (Algorithm \ref{alg:sleep}). Semantic replay initializes CTX from random activity and allows recurrent cortical dynamics to recover learned semantic attractors. As in episodic replay, sparsity is increased to isolate patterns corresponding to individual concepts (Fig. \ref{fig:phase_b}b). These patterns are then projected to MTL-semantic through initially random connections. The rationale mirrors that of episodic replay: increasing sparsity produces non-overlapping patterns that can be efficiently learned via competitive learning. As a result, receptive fields in MTL-semantic come to represent individual concepts identified by CTX (Fig. \ref{fig:phase_b}d). Consequently, while MTL-sensory presents a mixed selectivity to different concepts, MTL-semantic neurons show an absolute selectivity to these concepts (Fig. \ref{fig:phase_b}d, e), with mutual information becoming maximal rapidly after semantic replay is initiated.
\newline\newline
Overall, these results show that semantic replay implements semantic-to-episodic consolidation in the circuit model, as proposed in the cognitive model. This provides a mechanistic account of how learned semantic representations can be transferred from neocortex to MTL, giving rise to concept-like cells \shortcite{quiroga_2005,quiroga_2012} that are further examined below.
\begin{figure}
    \centering
    \includegraphics[width=\linewidth]{Figures/Figure_3.png}
\caption{\textbf{Semantic replay forms abstract representations in episodic memory.} \textbf{a}: Semantic-to-episodic consolidation via semantic replay. Cognitive model (left) and circuit implementation (right). During semantic replay, a random CTX pattern converges to a concept-specific attractor, which is then projected to MTL-semantic, creating semantic-specific receptive fields in MTL-semantic and enabling episodic traces to incorporate semantic information.  
\textbf{b}: Example neuronal activity during semantic replay. A random initial CTX state evolves into a complete pattern coding for a single concept, which drives activity in MTL-semantic. A combination of one-shot and slow learning again make neurons that are randomly slightly more sensitive to cortical patterns selective to cortical semantic representations. 
\textbf{c}: Final MTL-semantic$\leftarrow$CTX weight matrix, showing selective mapping from cortical semantic codes to MTL-semantic neurons.}
    \label{fig:phase_b}
\end{figure}
\begin{figure}[t]
  \ContinuedFloat
  \captionsetup{list=off} % don't add a new list-of-figures entry
  \caption{\textit{(continued)} \textbf{d}: Example activity during wake after semantic replay. MTL-semantic neurons now code selectively for concepts, complementing MTL-sensory representations.  
\textbf{e}: Selectivity index of neurons in MTL-sensory and MTL-semantic populations, showing that semantic replay induces sharper concept selectivity in MTL-semantic.  
\textbf{f}: Schematic of classifier training. Classifiers decode concept identity ($A_i$, $B_j$) from either MTL-sensory or MTL-semantic activity.  
\textbf{g--i}: Decoder performance across conditions. \textbf{g}: In-distribution decoding shows both MTL-sensory and MTL-semantic reliably code for concepts. \textbf{h}: Compositional out-of-distribution generalization. MTL-semantic neurons enable successful decoding of unseen concept combinations, while MTL-sensory neurons fail. \textbf{i}: Semantic out-of-distribution generalization. MTL-semantic neurons outperform MTL-sensory neurons in decoding novel semantic relationships. Statistical comparisons: *** $p < 0.001$, ns = not significant.}
\end{figure}
\newpage
\subsection*{Variable sparsity during replay controls the complexity of the vocabulary learned}
So far, episodic replay has been constrained to high sparsity, a regime that decorrelates episodic subpatterns and promotes the extraction of \textit{simple concepts} (e.g. tomato or cheese). However, if replay is allowed to recruit a larger number of active neurons (decreasing sparsity), the network can in principle recover patterns corresponding to higher-order combinations of concepts. For example, replay may jointly recover semantic components such as \textit{crust}, \textit{cheese}, and \textit{tomato} by allowing a less sparse pattern to emerge (Fig. \ref{fig:higher-order}a).
\newline\newline
To test this idea, we extended the replay dynamics to allow variable sparsity. We define the \textit{semantic load} as the ratio between the number of active neurons during replay and the number active during wakeful episodic encoding. Intuitively, semantic load measures how many concepts are simultaneously expressed in a replayed pattern. In the current setting, replay events have a semantic load of either 1 or 2 (Fig. \ref{fig:higher-order}b), such that MTL attractor dynamics recover a pattern associated  with either single concept (e.g. $A_1$) or a conjunction of two concepts (such as $A_1B_2$), depending on how many neurons are allowed to remain active. The resulting MTL activity is then projected to CTX, where a separate cortical subregion learns these higher-order combinations.
\newline\newline
Under this regime, replay with semantic load 2 induces in CTX a competition between patterns corresponding to specific combinations $(A_i,B_j)$, that is, complete episodic conjunctions rather than isolated concepts. Over time, CTX develops selective receptive fields for these combinations (Fig. \ref{fig:higher-order}d). As a result, cortical activity during wake contains neurons tuned to each $A_i$, to each $B_j$, and to specific conjunctions $(A_i,B_j)$ (Fig. \ref{fig:higher-order}c). In this way, cortical dictionary learning extends from simple cortical words (\textit{atoms} represent simple concepts) to higher-order compositional words (\textit{atoms} represent complex concepts).
\newline\newline
When examining CTX neurons representing complex concepts, we found that learning reorganizes activity into a conjunctive code, with each timestep representing an abstracted individual episode (Fig. \ref{fig:higher-order}d). Consistently, the mutual information between these neurons and complex latent factors of the form $A_iB_j$ increased sharply once variable sparsity was introduced (Fig. \ref{fig:higher-order}e). By contrast, the normalized mutual information of CTX simple-concept neurons and MTL-sensory activity with these complex latent factors remained notably below maximal (Fig. \ref{fig:higher-order}e). Finally, the formation of new receptive fields was favoured by an increased likelihood of the corresponding episode during wakefulness (Fig. \ref{fig:higher-order}f).
\newline\newline
This transformation is also reflected in corticocortical connectivity. As higher-order cortical words form, CTX recurrent synapses encode the relational structure between simple concepts and their conjunctions (Fig. \ref{fig:higher-order}e). It should be noted this structure is asymmetric: a conjunction such as $A_1B_1$ strongly predicts its constituent $A_1$, whereas $A_1$ only weakly predicts the specific conjunction $A_1B_1$. The resulting connectivity matrix is therefore strongly asymmetric.
\newline\newline
Together, these results show that replay sparsity controls the level of abstraction learned by the cortical dictionary. Highly sparse replay favours isolated concepts, whereas less sparse replay allows CTX to consolidate higher-order cortical words corresponding to concept conjunctions. In the pizza example, replay may first extract individual ingredients such as \textit{cheese}, \textit{tomato}, or \textit{crust}, and later consolidate pizza as a combination of the three, caputring increasingly episode-like semantic structure.
\begin{figure}[h!]
    \centering
    \includegraphics[width=0.95\linewidth]{Figures/Figure_4.png}
\caption{\textbf{Learning higher-order concepts based on simpler representations is facilitated by the presence semantic representations during replay.}}
    \label{fig:higher-order}
\end{figure}
\begin{figure}[h!]
  \ContinuedFloat
  \captionsetup{list=off} % don't add a new list-of-figures entry
  \caption{\textit{(continued)} \textbf{A}: Cognitive model (left) and circuit implementation (right) of higher-order concept learning. Multiple sensory–semantic episodes are stored in MTL. Variable sparsity during replay (semantic load) allows the recovery of higher-order associations, which are then consolidated into CTX$\leftarrow$MTL and  CTX$\leftarrow$CTX connections.\textbf{B}: Example activity during replay with semantic load 1 (top) or 2 (bottom). When one concept is replayed, CTX develops selective coding for single concepts. When two concepts are replayed simultaneously, CTX develops receptive fields for their combinations, supporting higher-order abstractions.  
\textbf{C}: Example activity during wake after training. CTX contains neurons selective to individual concepts $A_i$ and $B_j$, as well as to combinations ($A_i, B_j$).  
\textbf{D}: Final CTX$\leftarrow$MTL weight matrix. Cortical neurons integrate inputs from both MTL-sensory and MTL-semantic at different levels of abstraction.
\textbf{E}: Final CTX$\leftarrow$CTX weight matrix, showing asymmetric recurrent connectivity consistent with conditional probabilities across concepts. Block-like and diagonal substructures capture both simple and compositional units.  
\textbf{F}: Causal role of semantic replay. Random permutation of MTL-semantic$\leftarrow$CTX connections during higher-order learning significantly reduces episode selectivity in CTX (*** $p<0.001$), demonstrating that semantic representations actively scaffold the emergence of compositional structure. Figs. \ref{fig:experimental_support_2}b-e are obtained for the network with best accuracy in Fig. \ref{fig:experimental_support_2}f, see Fig. \ref{fig:higher_order_supp} for results on the network with median accuracy (key results still hold).}
\end{figure}
\newpage
\subsection*{Normative benefits of semantic representations in episodic memory}
The emergence of semantic representations in the model's Medial Temporal Lobe allows us to test their contribution to episodic memory and subsequent consolidation. The aim of this section is to clarify what are the normative benefits of embedding semantic structure within episodic memory. This question is particularly important in light of a longstanding tension between theory and experiment. On the one hand, classical attractor theories of episodic memory predict optimal performance for orthogonalized, uncorrelated patterns. On the other hand, experimental data suggest that episodic representations are far from unstructured, and instead contain clear semantic organization shared across memories. Our model provides a framework to examine why such structure might nevertheless be beneficial, and what computational trade-offs it introduces for recall, replay, and consolidation.
\newline\newline
Throughout this section, we will consider 3 network modes (Fig. \ref{fig:normative}a): \textit{semantics present} (no modification, MTL-semantic contains semantic representations from CTX), \textit{semantics random} (activity of MTL-semantic contains random activity at each timestep), and \textit{semantics absent}  (MTL-semantic is ignored during MTL-attractor dynamics). 
\subsubsection*{Semantic representations in episodic memory induce a generalization bias in replayed and consolidated representations}
Following up on new theoretical studies proposing that replay optimizes generalization in Complementary Learning Sytems (GO-CLS, \citeA{sun_2023}), we hypothesize semantic representations can play a role in the generalization of rerpresentations consolidated via replay.
\newline\newline
First, we investigated whether the presence of semantic content during episodic encoding has an effect on the generalization of newly learned concepts. To do so, we pre-trained network that learned only simple concepts (by limiting semantic load during replay), and created 3 copies: \textit{semantics present}, \textit{semantics random}, and \textit{semantics absent} (Fig. \ref{fig:normative}a). Then, we increased the maximum number of neurons active during replay (semantic load), allowing it to consolidate complex concepts that were combinations on the previously learned simple concepts (e.g. $A_1B_2$ is a combination of $A_1$ and$B_2$). We then measured the accuracy of CTX complex representations at test time (generalization).  We observed that as Signal-to-Noise Ratio (SNR) decreases, there is a higher gap between the generalization of cortical representations between networks with present and absent (or random) semantic representations (Figs. \ref{fig:normative}b and \ref{fig:normative}c , shadow X). This gap persists for intermediate level of noise, until the accuracy of both converge to chance (which expected from the noise in the teaching signal). (Figs. \ref{fig:normative}b and \ref{fig:normative}c, shadow X).
\newline\newline
To understand the origin of this cortical generalization gap, we next examined the replayed representations themselves. Two slightly different mechanisms could be driving this effect: when an episodic pattern is replayed, MTL is projected to CTX, and the activated cortical neurons update their MTL receptive fields in the direction of the replayed medial temporal pattern. Therefore, the formation of cortical representations during replay depends on two factors: (i) that the \textit{right} cortical neurons are activated (and therefore updated), and (ii) the similarity of the replayed MTL-sensory pattern with semantically equivalent future inputs. We quantified both factors across networks by measuring (i) \textit{Cortical Semantic Clustering} (do semantically equivalent replayed patterns activate the same cortical neurons during replay?) and (ii) \textit{MTL-sensory Semantic Clustering} (are semantically equivalent awake MTL-sensory patterns more similar to the replayed MTL-sensory representation?).
\newline\newline
Both factors were positively influenced by the presence of semantic representations (Fig. \ref{fig:normative}c), with a qualitatively larger effect on Cortical Semantic Clustering. To determine which of these two factors better explains the generalization gap between networks with and without semantic representations, we next examined how each measure related to the generalization of the resulting cortical representations. We found that Cortical Semantic Clustering showed the stronger relationship with cortical generalization, suggesting that it constitutes the primary mechanism through which semantic structure biases consolidation toward more generalizable representations. This result is consistent with the idea that cortical representations are not determined by any single replay event, but instead emerge gradually through the accumulation of many replay-driven updates under a small learning rate. As a consequence, even if semantic structure also makes individual replayed MTL-sensory patterns less noisy, the dominant factor for cortical generalization is whether the \textit{appropriate} cortical neurons are consistently recruited across replay events. In other words, what matters most is not only that replay is cleaner, but that replay repeatedly drives learning onto the same cortical assemblies.
\newline\newline
One intuitive interpretation is that MTL-semantic provides a low-dimensional, error-correcting scaffold that remains tightly coordinated with cortex and helps constrain which cortical neurons are activated during replay. Under this view, when an episodic pattern is replayed in MTL, the semantic component biases cortical recruitment toward assemblies corresponding to the underlying concept, effectively signaling that the current noisy episodic trace should be consolidated as another instance of a previously learned semantic structure. Indeed, ignoring the CTX$\leftarrow$MTL-semantic pathway during replay consolidation eliminated the generalization gap (Fig. SX).
\newline\newline
These results suggest that semantic representations promote cortical generalization not only by denoising replayed medial temporal patterns, but more importantly by biasing replay toward the repeated recruitment of the appropriate cortical assemblies. This effect becomes especially pronounced for more complex, episodic-like representations, where semantic structure provides an organizing scaffold across highly overlapping episodic patterns, which can otherwise follow unstable pattern completion dynamics. Together, our findings support the view that semantics shape replay into a teaching signal that is inherently more suitable for forming abstract, robust, and generalizable cortical representations.
\subsubsection*{Semantic representations in episodic memory increase memory capacity when episodes are semantically in-distribution}
Next, we examined the effect of semantic representations in episodic recall. As a baseline, we first asked whether the medial temporal recurrent network displayed a generic increase in capacity, independently of the structure of the stored patterns. To do so, we compared recall performance for purely random sparse patterns against a classic sparse Hopfield network of matched size and sparsity (Fig. \ref{fig:capacity}a). In this setting, our network showed a minor capacity advantage for a small number of patterns, but a steepest cliff (catastrophic forgetting). This confirmed the presence of the standard memory oveload characteristic of episodic storage of random patterns.
\newline\newline
We then examined how recall evolves with an increasing number of stored patterns when these are sampled from our standard episode distribution (Fig. \ref{fig:normative}d), for different levels of SNR in the sampled patterns. One most noticeable striking difference common across all networks and noise levels is the presence of a U-shaped learning curve. This is expected given that the underlying distribution is fixed, and the recurrent connections can end up averaging out noise yielding a better performance. An alternative view is that in practice the information (entropy) carried per pattern is smaller for structured patterns (Fig. SX), eroding less the network capacity when being added.
\newline\newline
However, the dip is more pronounced for \textit{semantics absent} and \textit{semantics random} networks, with a bigger drop in recall for intermediate numbers of patterns stored than \textit{semantics present}. Comparatively, the difference between \textit{semantics present} and the rest increases with noise. By compressing episodes into MTL-semantic, the information per full MTL pattern does not increase, but the number of available neurons doubles. Therefore, the neurons in MTL-semantic can error-correct MTL-sensory representations (increasing MTL-sensory recall).
\newline\newline
Finally, when looking at the recall of \textit{semantics present} within MTL-semantic is notably higher than that of MTL-sensory (within \textit{semantics present}), when again looking at the ratio of information across MTL-semantic patterns and its available number of neurons. This offers an explanation to why the semantic content of memories is more easily recalled than its sensory counterpart, where distortions usually happen.
\newline\newline
Closing paragraph
\begin{figure}
    \centering
    \includegraphics[width=0.95\linewidth]{Figures/Figure_5.png}
    \caption{Caption}
    \label{fig:normative}
\end{figure}
\subsection*{The model explains different behavioural and circuit experimental results}
\subsubsection*{Behaviour: benefits of blocked learning, better recall of semantic information, MTL lesion effects in semantic memory}
\subsubsection*{Circuits: concept cells, increased representational overlap with consolidation, and synaptic engrams in fear conditioning}
Taken together, the previous results suggest that semantic representations in episodic memory can help compressing complex sensory episodes into schematic representations, facilitating both episodic memory and further semantic integration of such episodes. This potential behavioural advantage rises the hypothesis that learning across different components of the circuit model can be mapped to neural and synaptic changes observed in experimental settings.
\newline\newline
For example, recent work has tracked engram-specific changes in synapses within and across regions involved in Contextual Fear Conditioning (CFC, Figure \ref{fig:experimental_final}a) \shortcite{lee_2023}. These experiments focused on the medial prefrontal cortex (mPFC) and basolateral amygdala (BLA), which correspond to the neocortex (CTX) and the medial temporal lobe (MTL) in our model (respectively, Figure \ref{fig:experimental_final}a). In order to recapitulate some of the main experimental results, we frame our previous input $(A_i, B_j)$ in a Contextual Fear Conditioning paradigm. We therefore assume the different $A_i$ to encode sensory representations of different environmental components (for example, $A_1$ can be thought of as the Conditioned Stimulus (CS) during fear conditioning, and $A_2$ as the home cage in which the animals are kept (Figures \ref{fig:experimental_final}a, b). Similarly, $B_j$ represents internal valence states such as fear (Unconditioned Stimulus (US), $B_1$, Figure \ref{fig:experimental_final}b), with $B_2, B_3, ...$ allowing representations of hunger, excitement, arousal... We assume a pre-existing semantic representation of fear, which is induced by presenting multiple episodes with persistent activation of $B_1$ fear MTL-sensory cells (Figure SX). This allows us to identify CTX \text{fear} cells, which we later use to model freezing behavior (these cells were not identified in experiments). Following this, we present the network, within a single day, first with only the Conditioned Stimulus (Figure \ref{fig:experimental_final}b), and then with also the US (fear, modelling the effect of the electric shock).
\newline\newline
As in the previous section, the model can have variable sparsity levels during sleep. As expected, this results in separately replaying both only Context A (CS), and the joint episode (US, CS). In this sense, two orthogonal representations form in CTX: one encoding only the CS, and one encoding the full episode. Instead of both populations overlapping, the semantic relationship with one another is encoded in synaptic connections: this allows keeping separate representations of context (CS), fear (US), and the episode, while storing the fact that the episode involved both, as well as being able to recollect the episode from partial cues (Supp. Fig. X).
\newline\newline
In experiments, a subpopulation of PFC and BLA neurons were tagged as \textit{engram} (active during episode encoding) or \textit{recall} (active during episodic recall). Then, photostimulation of different circuits allowed for direct measurement of synaptic strength across the tagged subpopulations, as shown in Fig. \ref{fig:experimental_final}c. Results indicated that, after consolidation, engram cells in PFC were more strongly connected to other engram cells in PFC than non-engram cells (not active during encoding, Fig. \ref{fig:experimental_final}c1), which in our model is explained by initial recurrent strengthening between newly following newly consolidated concepts. Similarly, connections projecting PFC engram cells to BLA engram cells were stronger than those connected to BLA neurons not active during encoding (Fig. \ref{fig:experimental_final}c2). Finally, connections from engram neurons to \textit{only recall} neurons (not active during encoding, but active during recall) were significantly stronger than connections to non-engram neurons, which was again reproduced by our model (Fig. \ref{fig:experimental_final}c3).
\newline\newline
Following fear conditioning, researchers further studied behavioural and circuit changes related to extinction (continued presentation of the Conditioned Simulus in the absence of the US (electric shock) (Fig. \ref{fig:experimental_final}d). Behavioural extinction of fear (Fig. \ref{fig:experimental_final}e) was correlated with a decrease in the significance of the synaptic strnegth difference across engram/non-engram subpopulations (Fig. \ref{fig:experimental_final}f). In the model, this behavioural extinction (Fig. \ref{fig:experimental_final}e) also results in a disruption of the original engram synaptic structure (Fig. \ref{fig:experimental_final}f). Although these experiments did not examine for conditioning after extinction, we further validated wether the model could explain the notion of \textit{savings} (original memory becomes inaccessible rather than erased). To do so, we present the model with a UC + CS sensory input after extinction, and measure CTX-only (PFC-only) memory recall of the episode. As hypothesized, a network that underwent fear conditioning + extinction scores significantly higher than a naive network with no CFC or extinction (Fig. \ref{fig:experimental_final}g), indicating that the memory was dormant rather than erased during extinction. This result can be mechanistically explained taking into account the idea semantics being encoded in recurrent cortical synapses as conditional probabilities. During extinction, the probability of the fear episode following the conditioned stimulus drifts from one (the first time) to 0 (never). However, both fear and context have maximal likelihood given the episode, maintaining those synapses intact (Fig. SX), and allowing recall in the case a CFC episode is experienced again.
\begin{figure}
    \centering
    \includegraphics[width=0.9\linewidth]{Figures/Figure_8.png}
    \caption{Caption}
    \label{fig:experimental_final}
\end{figure}
\clearpage
\begin{comment}
\subsection*{Episodic replay integrates sensory-semantic information into semantic memory: semantic refinement of cortical engrams}
During phase C of learning, the circuit is probed to refine cortical representations with the semantic input received from MTL-semantic (Fig. \ref{fig:phase_c}a). In this sense, neurons in CTX are expected to expand their receptive fields, now receiving meaningful projections from both MTL-sensory and MTL-semantic. In order to allow conjunctions of sensory and semantic patterns to be replayed together, both sub-regions in MTL share recurrent connectivity, storing the overall MTL activity with Hebbian plasticity. Similarly to phase A, the replayed activity corresponds to sub-patterns of stored episodes, but note how now within this sparser sub-patterns there is both a sensory and a semantic component. Mechanistically, phase C is equivalent to phase A, but now makes use of the joint sensory-semantic representations present in MTL.
\newline\newline
To test the degree in which this happens in the circuit, we first examine the MTL activity replayed during sleep. As in previous sections, sparsity during sleep is higher than during wakefulness, which allows recovering sensory-semantic sub-patterns (Fig. \ref{fig:phase_c}b). Note how, due to no distinction at the connectivity level (all MTL is a single recurrent lawyer) sensory and semantic patterns are coupled. 
\newline\newline
Next, we examine the evolution of the CTX$\leftarrow$MTL connectivity. During learning phase A, selective connections from MTL-sensory to CTX where already developed (Fig. \ref{fig:phase_a}f). Given that sensory and semantic activity is coupled during replay, now MTL-sensory drives activity in cortical neurons, which passively become tuned to MTL-semantic (Fig. \ref{fig:phase_c}c).
\newline\newline
While this last learning phase might seem redundant, in the next section we will see how it is crucial for cortical neurons to be able to read from MTL-semantic. In more realistic setups with lower selectivity in MTL-sensory, recall is more noisy in MTL-sensory than in MTL-semantic. This would make retrieval of previous experiences harder when only sensory information can be read from MTL.
\begin{figure}
    \centering
    \includegraphics[width=\linewidth]{Figures/Figure_4.png}
    \caption{Caption}
    \label{fig:phase_c}
\end{figure}
\end{comment}
\section*{Discussion}
In this study, we introduce a cognitive strategy by which episodic memory is improved through leveraging previously extracted semantic representations. Our model posits that, following an episodic-to-semantic consolidation process, semantic-to-episodic memory transfer reshapes the code of episodic memory itself. We further propose a biologically plausible neural circuit that implements such a cognitive model, successfully capturing the associated learning algorithm. Our circuit model aligns with and explains experimental data, and highlights the behavioural advantages of such a bidirectional learning mechanism.
\newline\newline
First, we demonstrated how episodic replay from hippocampal and para-hippocampal (medial temporal lobe, MTL) regions effectively supports semantic consolidation in cortical (CTX) networks. While this initial step of consolidation corresponds to the most widely studied phase of learning across subcortical-cortical interactions \shortcite{alvarez_1994, mcclelland_1995, oreilly_2014, sun_2023, guerreiro_2024, spens_2024} (to name just a few), our model exhibits key differences from previous models. To start, most models assume a faithful replay of awake subcortical and cortical activity, with abstraction happening directly in CTC$\leftarrow$CTX (neocortex to neocortex) connections. This is problematic for two reasons: first, the assumption that subcortical replay can reinstate perfectly previous cortical activity is not given \shortcite{swanson_2020}. Secondly, this relies on meaningful (to some degree) representations already existing in neocortex, when it is precisely via semantic consolidation that these representations should be created in the first place. For example, in \shortciteA{chrysanthidis_2022}, semantics are also extracted by forming block-like structure of recurrently connected neurons, but the high specificity of these neurons is already given. Furthermore, many studies assume a copy-paste mechanism between cortex and hippocampus, either directly via one-to-one connections \shortcite{alvarez_1994, tome_2022, guerreiro_2024, spens_2024} or perfect pattern-separating random connections \shortcite{sun_2023}, with less emphasis placed on how synapses mapping one to the other can be formed with biologically plausible learning mechanisms (although previous effort has been made in the context of backpropagation-like error-driven learning rules \shortcite{singh_2022}). In contrast, here we have proposed that increasing sparsity levels during replay can lead to CTX$\leftarrow$MTL synapses becoming selective to frequent episodic sub-patterns. This effectively decorrelates the training input provided to these connections, which could otherwise lead to representational collapse using standard Hebbian mechanisms. 
\newline\newline
Next, we proposed a mechanism whereby the semantic structure imprinted in CTX$\leftarrow$CTX connections can be transferred into the coding scheme of MTL. Both the notion that semantics are extracted in CTX \shortcite{mcclelland_1995} and that the MTL encodes a latent representation of the environment \shortcite{quiroga_2005, aronov_2017, behrens_2018} have been present in the literature for a long time, but a mechanistic account of how these representations emerge in MTL remains unknown. We put forward the hypothesis that, after extracted in CTX, spontaneous cortical activity could transfer information encoded in cortical synapses to MTL$\leftarrow$CTX synapses, inducing a semantic code in MTL. We called this process semantic-to-episodic consolidation, and conceptualized it as a mirror of episodic-to-semantic consolidation. As such, instead of transferring episodic information contained in recurrent MTL connections to CTX representations, it transfers semantic information present in recurrent CTX connections to MTL representations. The idea that cortex-to-hippocampus teaching is an underexplored topic in complementary learning systems literature was recently highlighted in \shortciteA{liu_2024}, and implicitly explored in \shortciteA{spens_2024}. In the latter study, the authors explored the possibility of using a Hopfield network to store pairs of sensory-latent representations for episodic recall, in a similar manner that our MTL-sensory and MTL-semantic neurons form joint attractors in MTL recurrent connections. Again, however, the process by which these joint representations were formed was not fully specified, with the model instead employing a copy-paste strategy to generate them. A model that also studied the formation of joint sensory-semantic representations as attractors in hippocampus is the Tolman Eichenbaum Machine \shortcite{whittington_2020}. In this study, the information in the sensory input was already highly semantized (i.e. using a one-hot encoding of sensory features), and the abstract graph position of the stimuli also given a priori. The mechanism for obtaining these representations in the first place was also not explored. In addition to establishing a mechanism for concept cell formation in MTL, our results show that this form of representations is more suitable for rapid learning than mixed representations, both in In-Distribution and Compositional Out-Of-Distribution environmental statistics. This allows MTL to compositionally encode never-seen episodes as long as these correspond to rearrangements of previously seen semantic representations. This is in agreement to previous results in continual learning in artificial neural networks \shortcite{chaudhry_2020}. We also showed how maintaining sensory low-level representations of the environment is crucial for encoding episodes that are Semantic Out-Of-Distribution (that is, when new episodes cannot be encoded using previous semantic representations). 
\newline\newline
Our model has shown how even in non-trivial, highly overlapping scenarios like a random projection of the original input one can obtain concept cells in MTL. Moreover, the model offers a mechanism via which the formation of these representations is favored by blocked (as opposed to interleaved) learning, as demonstrated in humans \shortcite{flesch_2018}. In our model, blocked learning provides the MTL with many different examples of episodes containing the same concept, which are averaged out in its recurrent connections, facilitating the replay of its corresponding prototypical representation \shortcite{posner_1968}. The model also explains why recall of episodes with a higher semantic component is better than that of episodes that, from the perspective of the learned semantic structure, appear random \shortcite{lin_2021}. In classic models of recall based on attractor models, overlaps between episodes is seen as detrimental, reducing the network capacity \shortcite{hopfield_1982, treves_1994}. However, semantic representations have a structured overlap, where the overlap is not stochastic but instead corresponds to the two episodes containing semantic overlaps. Thus, rather than reducing recall, this overlap acts as an error-correcting code \shortcite{albesagonzalez_2025}. Moreover, semantic representations not only show better recall, but also actively improve the recall of the sensory component. Similarly, semantic representations induce replaying activity patterns that resemble more activity during wakefulness. These results can help explain why learning schema-congruent information is much faster than that of completely new episodes \shortcite{tse_2007}. While many models have focused on explaining how faster neocortical learning could happen at the cortex level \shortcite{mcclelland_2013, guerreiro_2024}, our results suggest that this could be in part because the replayed activity in the first place is less noisy. Understanding the mechanisms of positive transfer learning is one of the cornerstones of modern artificial intelligence, where inducing models to use previously learned information to learn new tasks is still a challenge \shortcite{iman_2023}.
\newline\newline
Following this potential link between replay involving semantic representations and continual learning, we showed that closing the loop (performing episodic-to-semantic consolidation of episodes that contain both sensory and semantic representations) allows the formation of higher-order compositional representations in CTX \shortcite{kurth_2023}, in turn inducing a complete semantic description of the causal links across concepts in different levels of abstraction \shortcite{barlow_1972}. Not only higher-order representations (for example pizza is a conjunction of the different ingredients it contains) were developed, but CTX$\leftarrow$CTX connections captured the semantic relationships across these different hierarchical levels (such as the fact that pizza is much more related to cheese than cheese is to pizza). Crucially, these forms of asymmetric relationships cannot be captured by classic models of Hebbian learning, as recently shown in \shortciteA{albesagonzalez_2025}. Finally, our framework might be related to the phenomenon of infantile amnesia \shortcite{alberini_2017}. In the model, the semantic load quantifies how many distinct concepts are bound within a neural representation, ranging from low (highly conceptual) to high (highly episodic) values. Furthermore, we showed that consolidating representations with a high semantic load (those combining many concepts) requires pre-existing semantic structure to scaffold them. It could be that early in development, when semantic representations are still forming, this scaffolding is absent, making the consolidation of rich, episodic memories difficult. Consequently, infants may initially consolidate only simple, concept-like representations (e.g., family members or locations), while conjunctive memories combining these elements fail to stabilize. As semantic representations mature, episodes involving multiple known concepts can finally be consolidated. It would be interesting to study whether this coincides with the end of infantile amnesia.
\newline\newline
The model faces several limitations at different levels. For example, the model omits hash-like, episode-specific representations in episodic memory (as have been observed experimentally \shortcite{kolibius_2023, chettih_2024} and included in previous models \shortcite{sun_2023, chandra_2025}). An open question in neuroscience is how episodic representations solve the trade-off between perfect pattern separation and using overlapping patterns, more aligned with the hippocampus as a cognitive map that includes behaviourally-relevant latent representations. Further studies using a combination of shared and non-shared neuronal subspaces across similar episodes could explore a combination of both strategies might facilitate learning and optimal behaviour (efforts in this direction have been made before \shortcite{schapiro_2017}). Additionally, while we have slightly moved away from completely orthogonal input representations usually used in biologically-plausible mechanistic models of consolidation, the input used here was still relatively simple, with more complex representations still requiring gradient-based methods. How the brain might use local plasticity rules to learn from more realistic input remains an open problem.
\newline\newline
The results presented here could finally be of significance for modern artificial neural networks. In the recent years, many authors have highlighted episodic memory as a missing piece in many Artificial Intelligence (AI) paradigms \shortcite{fountas_2025, dong_2025, pink_2025}.  Our model suggests that episodic memory should not simply be added as a fast external buffer, but rather be continuously reshaped by semantic knowledge. This bidirectional interaction yields two advantages directly relevant for AI: (i) more robust recall of past experiences through semantic “error correction,” and (ii) the ability to rapidly learn new combinations of previously known concepts, supporting compositional generalization. Such designs may allow modern neural networks to leverage prior structure in order to accelerate continual learning and adapt flexibly to new environments.
\newline\newline
Altogether, our study shows that bidirectional consolidation can be captured by a biologically plausible circuit model relying only on local plasticity rules. We demonstrated how sparsity during replay enables cortical extraction of compositional structure, how semantic feedback supports the emergence of concept cells in MTL, and how the presence of semantic codes enhances recall, replay quality, and higher-order concept learning. These results extend standard consolidation models by providing a mechanistic explanation for the formation of semantic representations in MTL and their role in scaffolding both new episodic and semantic memories.
\newpage
\section*{Methods}\label{methods}
\subsection*{Episode Generation Protocol}
The Episode Generation Protocol (EGP) is a generative process that samples neuronal activity in the sensory layer. Our proposed EGP follows a similar approach to that described in \shortciteA{albesagonzalez_2025}, but is outlined here again for completion:
\subsection*{Episodes, episode attributes and episode concepts}
All \textit{episodes} share a structure, such that every episode contains one \textit{episode concept} per \textit{episode attribute}. This means that, given episode attributes $A$ and $B$, with each episode attribute representing a set of episode concepts $a\in A, b\in B$, then the set of all possible episodes is defined as:
\begin{equation}
    E = \left\{ \{a, b\} \mid a \in A, b \in B \right\}
\end{equation}
In other words, an episode $e\in E$ is a collection of episode concepts $a$, $b$ such that element $a$ belongs to episode attribute  $A$ and element $b$ belongs to episode attribute $B$. In this study we have fixed $a$ to have 5 possible values $a\in \{A_1, A_2, A_3, A_4, A_5\}$ and likewise for $b\in \{B_1, B_2, B_3, B_4, B_5\}$.
\newline\newline
\subsection*{Episode Probability Distribution}
As a generative process, our EGP samples episodes as a previous step to sampling sensory input. We do this by fixing a probability distribution over episodes:
\begin{equation}
    1 = \sum_{A_i, B_j}P(e = \{A_i, B_j\}) 
\end{equation}
such that $P(e = \{A_i, B_j\}) \equiv P(i,j)$ imposes the likelihood of an episode with latents $(A_i, B_j)$ to be sampled when generating episodes.
\subsection*{Episode-Input Mapping}
Ultimately, the EGP samples sensory layer activity. For this reason, one also has to define how a sampled episode is mapped into sensory input $X^\textrm{SEN}$. Every concept $c$ in $A\cup B$ has a set of associated neurons $\textrm{SEN}_{c}$ in $X^\textrm{SEN}$ such that, given an episode $e$
\begin{equation}
     X^\textrm{SEN}_{i} = \begin{cases} 
      1 & \textrm{if} \;\; i \in \textrm{SEN}_{c} \;\;\textrm{and}\;\; c \in e\\ 
      0 & \textrm{else} \\
   \end{cases}
\end{equation}
In order to define $\textrm{SEN}_c$, we specify the number of attributes $N_\textrm{attributes}$, in this case 2 ($A$ and $B$). Then, we define the number of concepts per attribute ($|A|$ and $|B|$, here 5 and 5, from this follows the fact that $A_{i \in [1, 5]}$ and $B_{j \in [1, 5]}$). Then, we fix the number of neurons per concept $N_\textrm{concept}^A$ and $N_\textrm{concept}^B$, and define $\textrm{SEN}_{c}$ as:
\begin{equation}
    \textrm{SEN}_{c} = \{i : \textrm{start}(c) \leq i < \textrm{stop}(c)\}
\end{equation}
with $\textrm{start}(c) = (i - 1) \cdot N_\textrm{concept}^A$ if $c = A_i \in A$ and $\textrm{start}(c) =|A|\cdot N_\textrm{concept}^A + (j - 1) \cdot N_\textrm{concept}^B$ if $c = B_j \in B$, and similarly $\textrm{stop}(c) = i \cdot N_\textrm{concept}^A$ if $c = A_i \in A$ and $\textrm{stop}(c) =|A|\cdot N_\textrm{concept}^A + j \cdot N_\textrm{concept}^B$ if $c = B_j \in B$ (i.e. we generate a list of one-hot encodings for each concept $c$).
\newline\newline
To account for variability between different presentations of the same episode (not all episodes, even though they contain the same concepts, will be exactly the same), we further add some stochasticity by randomly picking $N_\textrm{swap}$ inactive neurons and $N_\textrm{swap}$ active neurons and inverting their activity (a total of $2N_\textrm{swap}$ neurons randomly change their activity). This ensures that activity sparsity in the sensory layer is maintained (the number of swaps from 0 to 1 is the same as the number of swaps from 1 to 0).
\subsubsection*{Input temporal structure}
Generated input $X_\textrm{input}$ has dimensions $(N_\textrm{days}, T_\textrm{day})$. Every day, a total of number $T_\textrm{day}/T_\textrm{episode}$ episodes (of $T_\textrm{episode}$ length) are sampled, and within each timestep an independent noise realization is sampled. 
\subsubsection*{Semantic Structure of an EGP}
Following \shortciteA{albesagonzalez_2025}, we refer to  \textit{semantic field theory} \shortcite{bussmann_2006} in order to define what is the \textit{semantic structure} of an EGP. According to this theory, the meaning of a word is not isolated but dependent on its relation to the rest of the words. While our task is not one of language, we can use this same paradigm to define the \textit{meaning} of episode concepts. In this sense, the meaning (semantics) of our episode concepts depends on how they are related to the rest (also note this is how semantics are established in a language dictionary, where each concept is defined via its relationship with the rest of concepts). 
\newline\newline
Following this proposal, we use the conditional probabilities of being present in an episode between episode concepts as a proxy for the semantic structure of an EGP. In other words, extracting the semantics of an EGP is equivalent to extracting how likely is one episode concept $y$ to be present in an episode $e$ if an episode concept $x$ is also present. We define the Semantic Structure $\mathcal{SS}$ of an EGP with $N^\textrm{concepts}$ different concepts $c_i$, as as a $N^\textrm{concepts}\times N^\textrm{concepts}$ matrix of the form:
\begin{equation}
    \mathcal{SS}_{ij} = P(c_i \in e|c_j \in e) 
    \label{eq:semantic_structure}
\end{equation}
For example, the intuition that pizza has a stronger semantic connection to cheese than cheese to pizza is reflected in their conditional probabilities (it is very common to have cheese in a pizza, but much less to have pizza in any dish containing cheese).  
\subsection*{Circuit Model}
\subsubsection*{Architecture and nomenclature}
\paragraph*{Regions and neural activity}
The circuit model contains 3 main regions: sensory cortex (SEN), neocortex (CTX), and medial temporal lobe (MTL). In turn, MTL is split into two sub-regions: MTL-sensory and MTL-semantic. We use $X$ to denote the neural activity, which is written as $X^\textrm{region}$ (for example $X^\textrm{CTX}$) when the region is to be made explicit. The same follows for pre-activation (synaptic) input, which is denoted by $\hat{X}$ (or $\hat{X}^\textrm{region})$. The activity of neuron $i$ in a region is $X^\textrm{region}_i$. The number of neurons in a region is denoted by $N_\textrm{region}$ (e.g. $N_\textrm{CTX}$).
\paragraph{Connectivity}
Connectivity matrices connecting a region \textrm{pre} to a region \textrm{post} take the form $W^{\textrm{post}\leftarrow\textrm{pre}}$, and a synapse connecting neuron $j$ in region pre with neuron $i$ in region post is $W^{\textrm{post}\leftarrow\textrm{pre}}_{ij}$. For example, $W^{\textrm{CTX}\leftarrow\textrm{MTL}}_{13}$ represents the connection of neuron 3 in MTL with neuron 1 in CTX. Connectivity is classified as either feed-forward (connecting two different regions), which includes  $W^{\textrm{MTL-sensory}\leftarrow\textrm{SEN}}$, $W^{\textrm{CTX}\leftarrow\textrm{MTL}}$, and $W^{\textrm{MTL-semantic}\leftarrow\textrm{CTX}}$, or recurrent (connecting a region with itself), which includes $W^{\textrm{CTX}\leftarrow\textrm{CTX}}$, and $W^{\textrm{MTL}\leftarrow\textrm{MTL}}$. All connections are plastic except $W^{\textrm{MTL-sensory}\leftarrow\textrm{SEN}}$, which is fixed. Given a connectivity matrix $W^{\textrm{post}\leftarrow\textrm{pre}}$, we call receptive field of a postsynaptic neuron $i$ $\textrm{RF}^{\textrm{post}\leftarrow\textrm{pre}}_i$ the vector representing the connections from region pre to postsynaptic neuron $i$
\begin{equation}
    \textrm{RF}^{\textrm{post}\leftarrow\textrm{pre}}_i = \{{W^{\textrm{post}\leftarrow\textrm{pre}}}_{kj}\} \;\; : \;k = i 
\end{equation}
\paragraph{Synaptic Initialization}
\paragraph{Sensory cortex to MTL-sensory}
In Figs. \ref{fig:phase_a} and \ref{fig:phase_b}, $W^{\textrm{MTL-sensory}\leftarrow\textrm{SEN}}$ is initialized as the identity, mapping to MTL-sensory an exact copy of $X^\textrm{SEN}$. In the rest of Figures, the matrix is generated randomly with the condition that all receptive fields are of equal size:
\begin{equation}
    || RF_i ||_0 = |\{j \in \{1, ..., N_\textrm{SEN}\} : W^{\textrm{MTL-sensory}\leftarrow\textrm{SEN}}_{ij} \neq 0\}| = W^{\textrm{MTL-sensory}\leftarrow\textrm{SEN}}_\textrm{size}\;\; \forall i
\end{equation}
and sum:
\begin{equation}
    \sum_j\{RF_i^{\textrm{MTL-sensory}\leftarrow\textrm{SEN}}\}_j = \sum_jW^{\textrm{MTL-sensory}\leftarrow\textrm{SEN}}_{ij} = W^{\textrm{MTL-sensory}\leftarrow\textrm{SEN}}_\textrm{max-in}\;\; \forall i
\end{equation}
where $|| RF_i ||_0$ denotes the $l_0$ norm (which counts the number of non-zero entries), $|A|$ denotes the number of elements in set $A$, $W^{\textrm{MTL-sensory}\leftarrow\textrm{SEN}}_\textrm{size}$ is a network parameter determining the number of non-zero connections each receptive field has, and $W^{\textrm{MTL-sensory}\leftarrow\textrm{SEN}}_\textrm{max-in}$ another network parameter determining the maximum sum of incoming connections a postsynaptic neuron has. In order to achieve both conditions, we generate $W^{\textrm{MTL-sensory}\leftarrow\textrm{SEN}}$ by randomly sampling each receptive field $i$ independently, following:
\begin{equation}
    \{RF_i\}_j =
\begin{cases}
\frac{W^{\textrm{MTL-sensory}\leftarrow\textrm{SEN}}_\textrm{max-in}}{W^{\textrm{MTL-sensory}\leftarrow\textrm{SEN}}_\textrm{size}}, & \text{if } j \in S, \\
0, & \text{if } j \notin S,
\end{cases}
\end{equation}
with $S$ a random subset of the presynaptic entries:
\begin{equation} 
S \sim \text{Unif}\big\{ T \subseteq \{1,\dots,N_\textrm{SEN}\} : |T| = W^{\textrm{MTL-sensory}\leftarrow\textrm{SEN}}_\textrm{size} \big\} \Longrightarrow  |S| = W^{\textrm{MTL-sensory}\leftarrow\textrm{SEN}}_\textrm{size}, 
\end{equation}
\paragraph{Feed-Forward connections}
Feed-forward connections are sampled from a Gaussian distribution with mean 0 and standard deviation $W^{\textrm{post}\leftarrow\textrm{pre}}_\sigma$ as a network parameter.
\begin{equation}
    W^{\textrm{post}\leftarrow\textrm{pre}}_{ij} \sim \mathcal{N}(0, W^{\textrm{post}\leftarrow\textrm{pre}}_\sigma)
\end{equation}
\paragraph{Recurrent connections}
Recurrent connections are initialized at 0. MTL connections are reset every \textit{day} (see Algorithm \ref{alg:train})
\paragraph{Mature and immature neurons} Neurons in the model can be either \textit{mature} or \textit{immature}. The intuition is: do these neurons contain meaningful representations via input received from another region (mature) or are mostly non-selective and random (immature)? More details on the dynamics and plasticity of mature and immature neurons can be seen below, but to summarize: (i) immature neurons have higher excitability during feed-forward input processing (not during pattern completion) and (ii) immature neurons have a higher learning rate in their incoming feed-forward connections and its recurrent connections frozen. The intuition is that by having higher excitability, during receptive field formation in feed-forward synapses, immature neurons win the competition for postsynaptic activity only if the presynaptic pattern is too far from the receptive fields that have already formed. By doing this, if a presynaptic pattern is similar to previously formed representations, it slightly changes the already formed representations. If not, a set of immature neurons win the competition, develop receptive field due to higher plasticity, and then become mature. This allows one-shot formation of new postsynaptic representations of presynaptic input when an input pattern is too different from those presented to the network before, and then refining these representations with a slower learning rate. Recurrent connectivity in CTX, which reflects statistical regularities across the representations formed, is only developed between mature neurons that actually code for an environmental variable, to avoid recurrent connections representing spurious correlations. Maturity in a region is denoted by $\textrm{IM}^\textrm{region}$, where $\textrm{IM}^\textrm{region}_i = 1$ if neuron $i$ in the region is immature, and 0 otherwise.
\subsubsection*{Outgoing and incoming homeostasis across types of connections}
Below, we formalize two regimes of homeostatic plasticity: outgoing homeostasis and incoming homeostasis. Outgoing homeostasis guarantees that the sum of the outgoing connections from neuron $j$ in a connectivity matrix $W$ does not exceed a certain value $W_\textrm{max-out}$ ($\sum_iW_{ij} \leq W_\textrm{max-out} \forall j$). Similarly, incoming homeostasis guarantees a bound in incoming connections ($\sum_jW_{ij} \leq W_\textrm{max-in} \forall i$). Depending on the type of connectivity we choose:
\begin{itemize}
    \item outgoing homeostasis for CTX$\leftarrow$CTX: this makes connectivity matrix reflect semantic structure as defined in \eqref{eq:semantic_structure}, as connection from neuron coding for concept $i$ from neuron connecting for concept $j$ becomes proportional to the conditional probability of neuron $i$ firing given neuron $j$ fires \shortcite{albesagonzalez_2025}
    \item incoming homeostasis in feed-forward connectivity: in this homeostatic regime, connection from neuron $j$ to neuron $i$ becomes proportional to $p(j \; \textit{active} | i\; \textit{active})$ \shortcite{albesagonzalez_2025}. The \textit{trick} is that if an initial one-shot receptive field makes neuron $i$ selective to concept $c$, receptive fields approximate the average presynaptic activity pattern (prototype) given the concept c is present (neuron $i$ is active), thus reinforcing higher selectivity in noisier regimes.
    \item Homeostatic reset in MTL$\leftarrow$MTL connections. In order to account for symmetric plasticity kernels in hippocampus \shortcite{mishra_2016, li_2023}, we simply have Hebbian learning in MTL$\leftarrow$MTL. Homeostasis is obtained by resetting MTL recurrent connections to 0 after each sleep cycle.
\end{itemize}
\subsection*{Network Operations}
\subsubsection*{Neuronal activation}
Synaptic input is mapped to neural activity via a \textit{K-winners-take-all} mechanism, in which the $K$ neurons with the highest pre-activation input $\hat{X}$ set their activity to 1, and the rest to 0. This can be written as
\begin{equation}
    X = \textrm{top-}K\big(\hat{X}\big)
\end{equation}
In  order to determine $K$ in a particular region, first we define
\begin{equation}
    K^\textrm{region}_\textrm{concept} \leq K^\textrm{region}_\textrm{episode}
\end{equation}
 such that each concept is represented by $K^\textrm{region}_\textrm{concept}$ neurons, and a typical episode during wakefulness contains $K^\textrm{region}_\textrm{episode}$ active neurons. In this context, given $K$ active neurons in a region, we define the semantic load $S_L$ to be the number of concepts represented by $K$ neurons, such that
 \begin{equation}
     K = S_L \cdot K^\textrm{region}_\textrm{concept}
     \label{eq:semantic_load}
 \end{equation}
In other words, by fixing the semantic load, one fixes how many neurons in total are active, making $K$ higher if more concepts are being represented. The maximum number of concepts that can be encoded in an episode is thus:
 \begin{equation}
     N^\textrm{region}_\textrm{concept} = \frac{K^\textrm{region}_\textrm{episode}}{K^\textrm{region}_\textrm{concept}}
 \end{equation}
 Therefore, we can define \textsc{Activate} (Algorithm \ref{alg:activate}) as an algorithm that maps presynaptic input $\hat{X}$, given a semantic load $S_L$, into neural activation $X$ in the following way:
\begin{algorithm}[h!]
\caption{Neuronal Activation}\label{alg:activate}
\begin{algorithmic}[1]
\Function{Activate}{$\hat{X}$, $\textrm{region}$, $S_L$}
    \State Compute number of active neurons:
        \[
            K \gets S_L \times K^\textrm{region}_\textrm{concept}
        \]
    \State Select top-$K$ entries:
        \[
            X \gets \textrm{top-}K(\hat{X})
        \]
    \State \Return $X$
\EndFunction
\end{algorithmic}
\end{algorithm}

\subsubsection*{Feed-forward input processing}
Input processing is computed for a single pre and post region at a time, with pre-activation input $\hat{X}^\textrm{post}$ obtained via
\begin{equation}
    \hat{X}^\textrm{post} = W^{\textrm{post}\leftarrow\textrm{pre}} \cdot X^\textrm{pre} + b^\textrm{post}\textrm{IM}^\textrm{post}
\end{equation}
$b^\textrm{post}$ is the extra excitability that immature neurons have in region post.
\subsubsection*{Pattern completion (recurrent input processing)}
Given a recurrent network, pattern completion is typically defined as an iterative update on neural activity $X^\textrm{region}$ that is initiated in an initial state $X^\textrm{region}_0$:
\begin{equation}
    X^{\textrm{region}}_{k+1} = \textrm{top-}K(W^{\textrm{region}\leftarrow\textrm{region}}\cdot X^\textrm{region}_{k})
\end{equation}
with $W^{\textrm{region}\leftarrow\textrm{region}}$ the region recurrent connectivity matrix, $K$ dependent on the specified semantic load $S_L$, and $k$ a timestep operating on a smaller timescale,  than the overall timescale $t$. This means that within a timestep $t$, there are pattern\_complete\_iterations iterations on temporal index $k$. This is used to define \textsc{PatternComplete}  as Algorithm \ref{alg:pattern_complete}:
\begin{algorithm}[h!]
\caption{Pattern Completion}\label{alg:pattern_complete}
\begin{algorithmic}[1]
\Function{PatternComplete}{$X_0$, $S_L$,   \texttt{pattern\_complete\_iterations}}
    \State $X \gets X_0$
    \For{$\text{iteration} = 1$ to $\texttt{pattern\_complete\_iterations}$}
        \State $\hat{X} \gets W^{\textrm{region} \leftarrow \textrm{region}} \cdot X$
    \State $X \gets \textsc{Activate}\bigl(\hat{X}, \textrm{region}, S_L \bigr)$
    \EndFor
    \State \Return $X$
\EndFunction
\end{algorithmic}
\end{algorithm}
\subsubsection*{Hebbian Learning}
The connection between a presynaptic neuron $j$ and a postsynaptic neuron $i$ subject to Hebbian learning is updated as:
\begin{equation}
            \Delta W_{ij} = \lambda_{ij}^{\textrm{post} \leftarrow \textrm{pre}}
            X_i^\textrm{post} X_j^\textrm{pre}
\end{equation}
for a presynaptic region pre and a postsynaptic region post, where $\lambda_{ij}^{\textrm{post} \leftarrow \textrm{pre}}$ is the learning rate from neuron $j$ in region pre to neuron $i$ in region post.
\paragraph*{Maturity-dependent learning rate (feed-forward connections)}
In order to combine one-shot and statistical learning in feed-forward synapses (see \shortciteA{albesagonzalez_2025}), the learning rate of feed-forward (post $\neq$ pre) connections depends on the \textit{maturity} of the postsynaptic neuron, which is given by $\textrm{IM}$. Neurons start in an \textit{immature} state $\textrm{IM}_i = 1$, allowing them to form relatively selective initial receptive fields. Once this initial receptive field is formed, the learning rate is reduced to capture the statistical regularities of the presynaptic patterns driving the postsynaptic neuron (the neuron is \textit{mature}, $\textrm{IM}_i = 0$). To do this, we define a quick learning rate
\begin{equation}
        \lambda^{\textrm{post}\leftarrow\textrm{pre}}_\textrm{quick} = \frac{W^{\textrm{post}\leftarrow\textrm{pre}}_\textrm{max-in}}{K^\textrm{pre}} 
\end{equation}
where again $W^{\textrm{post}\leftarrow\textrm{pre}}_\textrm{max-in}$ is a network parameter determining the maximum sum of incoming connections the matrix $W^{\textrm{post}\leftarrow\textrm{pre}}$ has for every postsynaptic neuron ($\sum_jW^{\textrm{post}\leftarrow\textrm{pre}}_{ij} \leq W^{\textrm{post}\leftarrow\textrm{pre}}_\textrm{max-in}$). Then, at each Hebbian update, we have
\begin{equation}
    \lambda_{ij}^{\textrm{post} \leftarrow \textrm{pre}} = (1 - \textrm{IM}^\textrm{post}_i)\lambda^{\textrm{post} \leftarrow \textrm{pre}}_\textrm{slow} + \textrm{IM}^\textrm{post}_i\lambda^{\textrm{post}\leftarrow\textrm{pre}}_\textrm{quick}
\end{equation}
where $\lambda^{\textrm{post} \leftarrow \textrm{pre}}_\textrm{slow}$ is a network parameter indicating the slow learning rate (used throughout all simulations except for the first receptive field formed at each postsynaptic neuron). Note how for an initial (before maturity) driving presynaptic pattern $X^\textrm{pre}$ this guaranties that the initial receptive field formed is $X^\textrm{pre}$ scaled such that the sum of incoming connections is exactly equal to $W^{\textrm{post}\leftarrow\textrm{pre}}_\textrm{max-in}$, as the $K$-winners-share-all mechanism imposes exactly $K_\textrm{pre}$ active neurons. 
\paragraph*{Maturity-dependent learning rate (CTX recurrent connections)}
Recurrent connections are designed to capture conditional firing probabilities between neurons. When individual neurons become mature and tuned to individual concepts, this allows capturing the environment semantic structure. To avoid capturing the statistical structure of immature neurons, and based on the fact that learning in CTX is favoured by the previous formation of schemas \shortcite{tse_2007}, we follow an opposite recipe. Hebbian updates on recurrent connections are only effective if both pre- and postsynaptic neurons are mature:
\begin{equation}
    \lambda_{ij}^{\textrm{region} \leftarrow \textrm{region}} = (1 - \textrm{IM}^\textrm{post}_i)(1 - \textrm{IM}^\textrm{post}_j)\lambda^{\textrm{region}\leftarrow\textrm{region}}_0
\end{equation}
In the case of recurrent connections, there are no slow and quick learning rates, but the parameter $\lambda^{\textrm{post}\leftarrow\textrm{pre}}_0$ is used to distinguish it from the maturity-dependent learning rate matrix 
\paragraph*{Hebbian learning algorithm}
Altogether, this defines \textsc{Hebbian} (Algorithm \ref{alg:hebbian}) as the following update in $W^{\textrm{post} \leftarrow \textrm{pre}}$:
\begin{algorithm}
\caption{Hebbian Learning}\label{alg:hebbian}
\begin{algorithmic}[1]
\Function{Hebbian}{$\textrm{post}$, $\textrm{pre}$}
    \State Obtain learning rate matrix:
    \If{$\textrm{post} \neq \textrm{pre}$}
        \[
            \lambda_{ij}^{\textrm{post} \leftarrow \textrm{pre}}
            \gets
            (1-\textrm{IM}^\textrm{post}_i)\,\lambda^{\textrm{post} \leftarrow \textrm{pre}}_\textrm{slow}
            + \textrm{IM}^\textrm{post}_i\,\lambda^{\textrm{post}\leftarrow\textrm{pre}}_{\textrm{quick}}
        \]
    \ElsIf{$\textrm{post} = \textrm{pre} = \mathrm{CTX}$}
        \[
            \lambda_{ij}^{\textrm{post} \leftarrow \textrm{pre}} 
            \gets 
            (1 - \textrm{IM}^\textrm{post}_i)(1 - \textrm{IM}^\textrm{pre}_j)\,
            \lambda^{\textrm{post}\leftarrow\textrm{pre}}_\textrm{0}
        \]
    \Else
        \[
            \lambda_{ij}^{\textrm{post} \leftarrow \textrm{pre}} 
            \gets 
            \lambda^{\textrm{post}\leftarrow\textrm{pre}}_0
        \]
    \EndIf
    \State Compute Hebbian update with activity outer product:
        \[
            \Delta W \gets \lambda^{\textrm{post} \leftarrow \textrm{pre}} \odot 
            \Bigl(X^{\textrm{post}} \otimes X^{\textrm{pre}}\Bigr)
        \]
    \State Update weights:
        \[
            W^{\textrm{pre} \leftarrow \textrm{post}} 
            \gets 
            W^{\textrm{pre} \leftarrow \textrm{post}} + \Delta W
        \]
\EndFunction
\end{algorithmic}
\end{algorithm}
\subsubsection*{Homeostatic Plasticity}
We use multiplicative normalization that guarantees, for every weight matrix $W^{\textrm{post}\leftarrow\textrm{pre}}$, the total sum of incoming or outgoing connections is capped to a certain value. Following \shortciteA{albesagonzalez_2025}, we use \textit{incoming homeostasis} as:
\begin{equation}
W_{ij}^{k+1}=
\begin{cases}
 W_{ij}^{k}\cdot\frac{W_\textrm{max-in}}{\sum_{j}W^{k}_{ij}} & \text{if }\sum_{j}W^{k}_{ij} = W_\textrm{max}^\textrm{in} + \epsilon_i, \;\; \epsilon_i > 0 \\[1.2ex]
W^k_{ij}& \textrm{else}
\end{cases}
\end{equation}
where $W^k_{ij}$ denotes the value of synapse $ij$ before applying homeostasis and $W^{k+1}_{ij}$ the value after applying homeostasis (for readability, we simply write $W \equiv W^{\textrm{post}\leftarrow\textrm{pre}}$).
Similarly, we have \textit{outgoing homeostasis} to be defined by:
\begin{equation}
W_{ij}^{k+1} =
\begin{cases}
 \frac{W_\textrm{max-out}}{\sum_{i}W^k_{ij}} & \text{if }\sum_{j}W^k_{ij} = W_\textrm{max}^\textrm{out} + \epsilon_j, \;\; \epsilon_j > 0 \\[1.2ex]
0 & \textrm{else}
\end{cases}
\end{equation}
This allows, for a presynaptic region pre and a postsynaptic region post, defining \textsc{Homeostasis} as Algorithm \ref{alg:homeostasis}.
\subsection*{Learning fixed points}
The combination of \textit{K-winners-take-all} ($\textrm{top-}K$) activation function with Hebbian and homeostatic plasticity as defined here has been recently shown \shortcite{albesagonzalez_2025} to make synaptic connections converge to conditional firing probabilities. For outgoing homeostasis, the fixed points are:
  \begin{equation}
     W_{ij}^{\textrm{pre}\leftarrow\textrm{post}, \; \textrm{out}}(\infty) \equiv \frac{W_\textrm{max-out}^{\textrm{pre}\leftarrow\textrm{post}}}{K^\textrm{post}}p(X^\textrm{post}_i = 1| X^\textrm{pre}_j = 1)
 \end{equation}
 where $K^\textrm{post}$ corresponds to the number of active neurons during plasticity. Given that the conditional probability can be at most 1, we define in this scenario $W_\textrm{max}$ as $\frac{W_\textrm{max-out}^{\textrm{pre}\leftarrow\textrm{post}}}{K^\textrm{post}}$. Similarly, for incoming homeostasis, the fixed points are
   \begin{equation}
   W_{ij}^{\textrm{pre}\leftarrow\textrm{post}, \; \textrm{out}}(\infty) \equiv \frac{W_\textrm{max-in}^{\textrm{pre}\leftarrow\textrm{post}}}{K^\textrm{pre}}p(X^\textrm{pre}_j = 1| X^\textrm{post}_i = 1)
 \end{equation}
 with $W_\textrm{max}=\frac{W_\textrm{max-in}^{\textrm{pre}\leftarrow\textrm{post}}}{K^\textrm{pre}}$. These results have an important associated intuition, which is that if neurons code for individual concept (have a concept selectivity close to 1), synaptic connections reflect the associated Semantic Structure \shortcite{albesagonzalez_2025} under outgoing homeostasis (as happens in CTX). In MTL-semantic$\leftarrow$CTX and CTX$\leftarrow$MTL, which follow incoming homeostasis, when initial one-shot receptive fields induce a selectivity to a concept close to 1, feed-forward synapses slowly approximate to the associated prototype, where each synapse is becomes proportional to the contribution of each presynaptic neuron to the concept encoded in the postsynaptic neuron.
\begin{algorithm}
\caption{Homeostasis}\label{alg:homeostasis}
\begin{algorithmic}[1]
\Function{Homeostasis}{$\textrm{post}$, $\textrm{pre}$}
\State $W \gets W^{\textrm{post}\leftarrow\textrm{pre}}$
    \For{each post-synaptic unit $i$}
        \If{$\sum_j W[i,j] > W^{\textrm{post}\leftarrow\textrm{pre}}_\textrm{max-post}$}
            \State Normalize:
                \[
                W[i,:] \gets W[i,:] 
                \times
                \frac{W^\textrm{}}{\sum_j W[i,j]}
                \]
        \EndIf
    \EndFor
    \For{each pre-synaptic unit $j$}
        \If{$\sum_i W[i,j] > \text{max\_pre\_connectivity}$}
            \State Normalize:
                \[
                W[:,j] \gets W[:,j] 
                \times
                \frac{\text{max\_pre\_connectivity}}{\sum_i W[i,j]}
                \]
        \EndIf
    \EndFor
\EndFunction
\end{algorithmic}
\end{algorithm}
\subsubsection*{Replay}
We define \textrm{Replay} as a learning mechanism involving a pre and a post region, such that $W^{\textrm{post}\leftarrow\textrm{pre}}$ extract statistical structure from $W^{\textrm{pre}\leftarrow\textrm{pre}}$ connections. A replay operation takes as argument a specific semantic load $S_L$, which determines the conceptual complexity of the extracted regularities. To do this, we first pattern complete an initial presynaptic pattern $X^\textrm{pre}_0$, defined as a Gaussian random vector of the same size as the pre region:
\begin{equation}
     X^\textrm{pre}_{0} \sim \mathcal{N}(0, I_{N_\textrm{pre}})
\end{equation}
where $N_\textrm{pre}$ is the size of region pre and $\mathcal{N}(0, I_{N_\textrm{pre}})$ is a multivariate Gaussian distribution with 0 mean and identity covariance. This initial random pattern is pattern completed to obtain $X^\textrm{pre}$, and then projected to the post region, to give $X^\textrm{post}$, using in both operations the same fixed semantic load. Then, Hebbian and homeostatic plasticity operations as defined in Algorithms \ref{alg:hebbian} and \ref{alg:homeostasis} are used in the (post, pre) region pair. Replay is computed in simulations following Algorithm \ref{alg:replay}. In the text, we call replay from MTL to CTX \textit{Episodic Replay}, because it finds regularities across recurrent synapses that contain episodic information, and replay from CTX to MTL-semantic \textit{Semantic Replay}, because in this case recurrent synapses contain semantic information.
\begin{algorithm}[h!]
\caption{Replay}\label{alg:replay}
\begin{algorithmic}[1]
\Function{Replay}{post, pre, $S_L$, \texttt{pattern\_complete\_iterations}}
    \State Sample Gaussian random vector in pre:
        \[
            X_0^\textrm{pre} \gets \text{randn}(N_\textrm{pre})
        \]
    \State $X^\textrm{pre} \gets \textsc{PatternComplete}\bigl(X^\textrm{pre}_0, S_L, \texttt{pattern\_complete\_iterations}\bigr)$
    \State $X^\textrm{post} \gets W^{\textrm{post}\leftarrow\textrm{pre}}\cdot X^\textrm{pre} + b^\textrm{post}\textrm{IM}^\textrm{post}$
    \State \textsc{Hebbian}(post, pre)
    \State \textsc{Homeostasis}(post, pre)
    \State \Return $X$
\EndFunction
\end{algorithmic}
\end{algorithm}
\subsection*{Learning Algorithm}
For each day of experience, assuming an input $X_\textrm{input}$ of shape $(N_\textrm{days}, T_\textrm{day})$, we define the following learning algorithm \textsc{Train} (Algorithm \ref{alg:train}), where \textrm{Wake} is defined as in Algorithm \ref{alg:wake} and the \textrm{Sleep} as in Algorithm \ref{alg:sleep}. All operations used have been previously defined, with these algorithms showing the order in which they are performed. \textsc{duration\_phase\_A} represents the number of days in which only episodic replay takes place (Fig. \ref{fig:phase_a}), and  \textsc{duration\_phase\_B} the number of days in which only MTL-sensory synapses are used for episodic replay (but MTL-semantic representations have been formed already, Figs. \ref{fig:phase_b}, \ref{fig:experimental_support_1}, \ref{fig:experimental_support_2}, except Fig. \ref{fig:experimental_support_2}f). From day \textsc{duration\_phase\_B} on, MTL replays both sensory and semantic representations. Therefore, during \textsc{duration\_phase\_A} days, episodic memories containing only sensory information are stored in MTL during Wake (Algorithm \ref{alg:wake}). During Sleep (Algorithm \ref{alg:sleep}), cortical abstractions of sensory information are formed. After \textsc{duration\_phase\_A} days, cortical abstractions are transferred to MTL-semantic via semantic replay during Sleep. Finally, after \textsc{duration\_phase\_B} days, episodic memories contain both sensory and semantic information, and so does episodic replay. This allows higher-order abstraction, resulting in the formation of new semantic representations that are based on previous abstractions.
\begin{algorithm}[h!]
\caption{Learning Algorithm Per Day}\label{alg:train}
\begin{algorithmic}[1]
\Function{Train}{$X_\textrm{input}$}
    \For{$\textrm{day} = 1$ to $N_{\textrm{days}}$}
        \State Extract daily input:
            \[
            X^{\textrm{day}}_\textrm{input} \gets X_\textrm{input}[\textrm{day},:]
            \]
        \State \textsc{Wake}($X^{\textrm{day}}_\textrm{input}$)
        \State \textsc{Sleep}()
    \EndFor
\EndFunction
\end{algorithmic}
\end{algorithm}
\begin{algorithm}[h!]
\caption{Wake}\label{alg:wake}
\begin{algorithmic}[1]
\Function{Wake}{$X^{\textrm{day}}$}
        \State Reset MTL recurrent connections:
            \[
                W^{\textrm{MTL} \leftarrow \textrm{MTL}} \gets 0
            \]
    \For{$t = 1$ to $T_{\textrm{day}}$}
        \State $X^{\textrm{SEN}} \gets \textsc{Activation}\bigl(X^{\textrm{day}}[t], \textrm{SEN}\bigr)$
        \State $\hat{X}^{\textrm{MTL-sensory}} 
            \gets 
            W^{\textrm{MTL-sensory} \leftarrow \textrm{SEN}} 
            \cdot X^{\textrm{SEN}}$
        \State $X^{\textrm{MTL-sensory}} 
            \gets 
            \textsc{Activation}\bigl(\hat{X}^{\textrm{MTL-sensory}}, \textrm{MTL-sensory}\bigr)$
        \State $\hat{X}^{\textrm{MTL-semantic}} 
            \gets 
            \text{randn}(N_\textrm{MTL-semantic})$
        \State $X^{\textrm{MTL-semantic}} 
            \gets 
            \textsc{Activation}\bigl(\hat{X}^{\textrm{MTL\_semantic}}, \textrm{MTL\_semantic}\bigr)$
        \State $X^{\textrm{MTL}} 
            \gets 
            [X^{\textrm{MTL-sensory}}, \; X^{\textrm{MTL-semantic}}]$
        \State $\hat{X}^{\textrm{CTX}} 
            \gets 
            W^{\textrm{CTX} \leftarrow \textrm{MTL}} 
            \cdot X^{\textrm{MTL}}$
        \State $X^{\textrm{CTX}} 
            \gets 
            \textsc{Activation}\bigl(\hat{X}^{\textrm{CTX}}, \textrm{CTX}\bigr)$
        \If{day $\ge$ duration\_phase\_A}
            \State $\hat{X}^{\textrm{MTL-semantic}} 
                \gets 
                W^{\textrm{MTL-semantic} \leftarrow \textrm{CTX}} 
                \cdot X^{\textrm{CTX}}$
            \State $X^{\textrm{MTL-semantic}} 
                \gets 
                \textsc{Activation}\bigl(\hat{X}^{\textrm{MTL-semantic}}, \textrm{MTL-semantic}\bigr)$
            \State $X^{\textrm{MTL}}[N_\textrm{MTL-sensory}:] 
                \gets 
                X^{\textrm{MTL\_semantic}}$
            \State \textsc{Hebbian}$(X^{\textrm{MTL\_semantic}}, X^{\textrm{MTL\_semantic}})$
        \EndIf
        \If{day $\ge$ duration\_phase\_B}
            \State $\hat{X}^{\textrm{CTX}} 
                \gets 
                W^{\textrm{CTX} \leftarrow \textrm{MTL}} 
                \cdot X^{\textrm{MTL}}$
            \State $X^{\textrm{CTX}} 
                \gets 
                \textsc{Activation}\bigl(\hat{X}^{\textrm{CTX}}, \textrm{CTX}\bigr)$
        \EndIf
        \State \textsc{Hebbian}$(\textrm{MTL}, \textrm{MTL})$
        \State \textsc{Hebbian}$(\textrm{CTX}, \textrm{CTX})$
        \State \textsc{Homeostasis}$(\textrm{CTX}, \textrm{CTX})$
    \EndFor
    \State Increment day counter
\EndFunction
\end{algorithmic}
\end{algorithm}
\begin{algorithm}[h!]
\caption{Sleep}\label{alg:sleep}
\begin{algorithmic}[1]
\Function{Sleep}{}
\State Episodic Replay
    \For{$t = 1$ to $\texttt{sleep\_duration\_A}$}
        \If{day $\le$ duration\_phase\_B}
            \State \textsc{Replay}(\textrm{CTX}$, \textrm{MTL-sensory}$, $S_L= 1$, $\textrm{pattern\_complete\_iterations}=10$)
        \Else
        \State Sample $S_L$ randomly
        \State $S_L \gets \textsc{randint}(1, S_{L-\textrm{max}})$
            \State \textsc{Replay}(\textrm{CTX}$, \textrm{MTL}$, $S_L= S_L$, $\textrm{pattern\_complete\_iterations}=10$)
        \EndIf
    \EndFor

    \If{day $\ge$ duration\_phase\_A}
        \For{$t = 1$ to $\texttt{sleep\_duration\_B}$}
            \State \textsc{Replay}(\textrm{MTL-semantic}$, \textrm{CTX}$, $S_L= 1$, $\textrm{pattern\_complete\_iterations}=10$)
        \EndFor
    \EndIf
\EndFunction
\end{algorithmic}
\end{algorithm}
\clearpage
\subsection*{Data Analysis}
\subsubsection*{Neuronal Selectivity}
Given a vector $X$ that contains neuronal activity across T timesteps -with dimensions $(T, N_\textrm{neurons})$- and one that contains certain latent variables $Z$ -with dimensions  $(T, N_\textrm{latents})$, we do obtain the normalized values:
\begin{equation}
    X_\textrm{norm} = \frac{X - \langle X \rangle_{\textrm{time}}}{\sqrt{\langle X^2 \rangle_\textrm{time} - \langle X \rangle_\textrm{time}^2}} \;\; ; \;\; Z_\textrm{norm} = \frac{Z - \langle Z \rangle_{\textrm{time}}}{\sqrt{\langle Z^2 \rangle_\textrm{time} - \langle Z \rangle_\textrm{time}^2}}
\end{equation} 
And then obtain the matrix 
\begin{equation}
    S = \frac{1}{T}X_\textrm{norm}^TZ_\textrm{norm}
\end{equation}
which, for neuron $i$ and latent $j$ contains their correlation as:
\begin{equation}
    S(i, j) = \frac{1}{T}\sum_t X_{ti}Z_{tj}
\end{equation}
The Maximum Selectivity of neuron $i$ is the maximum selectivity value of neuron $i$ across all latents $j$:
\begin{equation}
    S_\textrm{max}(i) = \max_{j}{S(i, j)}
\end{equation}
\subsubsection*{Neuron Ordering}
Throughout this study, we re-order the neurons indices for visualization purposes. We do this with respect to a certain latent variable (for example the different \textit{concepts} $A_i$ and $B_i$ or the different episodes $A_iB_j$. This is done by imposing a threshold, and then for each possible discrete value the latent can take we obtain the neurons with a selectivity higher to that value with respect to the given latent. For example, assuming latents $A_1$ to $B_5$, the first neurons would be those with a selectivity higher than 0.9 to $A_1$, then the next neurons would be those selective to $A_2$, etc. All neurons that are not selective to any latent keep the original order at the end of the new order.
\subsection*{Quick Downstream Learning}
In Fig. \ref{fig:phase_b}, we test how well can a classifier learn from different representations. These representations are: MTL-sensory, MTL-semantic and \textit{labels}. While the first two are simply the simulated neural activity of the corresponding regions, \textit{labels} are simply lists one-hot encodings of the labels being predicted (i.e. simply the same as the output, repeated to match the size of MTL). This is to obtain a baseline of the most \textit{trivial} task possible, which would be mapping labels to labels. Then, for each panel in \ref{fig:phase_b}g, we initialize at random 100 classifiers consisting on a linear layer followed by a softmax activation. For each dataset $(X^\textrm{MTL-sensory}, Y^\textrm{latent})$, $(X^\textrm{MTL-semantic}, Y^\textrm{latent})$, and $(X^\textrm{labels}, Y^\textrm{latent})$ we train the same (initially) classifier to predict $Ai$ and $B_j$ that generated each representation $X^\textrm{region}$. We do this for only one epoch as a proxy for quick downstream learning facilitation. Each epoch consists of 100 days of experience, which amounts to 8000 samples.
\newline\newline
While obtaining the neural representations at each region, the model is frozen. This allows us to study 3 different regimes depending on whether the outside world distribution matches the original input distribution or not: (1) In-Distribution (input sampled from the same distribution) (2) Compositional Out-Of-Distribution (episodes had never been seen by the network, but the individual compositional blocks are maintained), and (3) Semantic Out-Of-Distribution (input generated using never-seen blocks, here implemented as a fixed random permutation of original input to maintain rest of statistics).  
\subsection*{Random projection of sensory inputs}
In Figs. \ref{fig:experimental_support_1} and \ref{fig:experimental_support_2} we forward the initially sampled episode through random connections. This connectivity is obtained by fixing a receptive field size $N_\textrm{RF}^\textrm{pre}$, then for each postsynaptic neuron $i$ randomly picking $N_\textrm{RF}^\textrm{pre}$ synapses $ij \in \textrm{RF}(i)$, and doing
\begin{equation}
    W^{\textrm{post}\leftarrow\textrm{pre}}_{ij} = W_\textrm{max}^\textrm{post}/N_\textrm{RF}^\textrm{pre}
\end{equation}
Thus, every postsynaptic neuron is randomly assigned $N_\textrm{RF}^\textrm{pre}$ presynaptic neurons, and has total incoming weights $W_\textrm{max}^\textrm{post}$. Given our top-k activation, this guarantees a much more distributed firing rate, which depends exclusively on the presynaptic pattern and not on the total sum of weights assigned at random.
\subsection*{Blocked vs interleaved training}
In order to model blocked training, we change input statistics during the first 100 days. For each concept $A_1$ to $B_5$, we have a 10-day block of 200 timesteps where every episode is sampled with equal probability as long as the corresponding concept is present. During the first block, all episodes with $A_i = A_1$ are equally likely (the rest have probability 0), then second block the same with $A_i = A_2$, until $B_j = B_5$. For interleaved training, we maintain original statistics but also increase the number of timesteps to 200 per day so the total amount of episodes is the same. 
\subsection*{Concept prototypes}
For a given neural representation $X^\textrm{region}$, and a concept $c$, one can define the prototype of $c$ in the region as
\begin{equation}
    \textrm{top-}K^\textrm{region}\big(\langle X^\textrm{region} \rangle_{c \in e}\big)
    \label{eq:prototype}
\end{equation}
that is, the activation of the average neural activity across all episodes in which concept $c$ is present. For example, the protototype in MTL-sensory of $A_1$ would be the average neural activity across episodes containing $A_1$. Note how this could also be extended to higher-order concepts (for example the prototype of $A_1B_3$).
\subsection*{Recall and replay}
In Fig. \ref{fig:experimental_support_2} we obtain two measures related to the episodic memories formed in MTL$\leftarrow$MTL connections: \textit{recall} and \textit{replay}. In both cases we have a training phase, where multiple episodes are sampled as usual, and all of which are stored in MTL recurrent connections following \textit{wake} (Algorithm \ref{alg:wake}). Then there is a test phase. For \textit{recall}, we (1) randomly sample the neural activity of MTL-sensory for one of the episodes presented during training, (2) impose that activity and leave MTL-semantic random, and finally obtain $X^\textmr{MTL}_\textrm{P.Completed}$ via the \textit{Pattern Complete} operation (Algorithm \ref{alg:pattern_complete}). Recall is defined as the cosine similarity between the encoded and the recovered representation. Where indicated in a legend as MTL-sensory or MTL-semantic, that means the cosine similarity measured only in that subregion (but the previeous procedure is maintained the same). To measure  replay, we simply follow Algorithm (\ref{alg:replay}), which means doing pattern completion on a random initial activity pattern. Then, we measure the cosine similarity of MTL-sensory with each of the different prototypes (Eq. \eqref{eq:prototype}), and obtain the highest score across all prototypes. Intuitively, because there was no original pattern presented, and because we are interested in measuring replay as an initial generalization step, we measure how well do replayed activity patterns reflect latent prototypes (as these patterns are the ones that will be used to form semantic representations in CTX).
\subsection*{Simulation parameters}
Default input, temporal, and model parameters can be found in Tables \ref{tab:default_input}-\ref{tab:default_model}. Figure-specific parameters can be found in Tables \ref{tab:figure_2}-\ref{tab:figure_6}.
\begin{table}[ht]
\centering
\caption{Summary of Episode Generation Protocol parameters}
\label{tab:default_input}
\begin{tabular}{lll}
\toprule
\textbf{Variable} & \textbf{Description} & \textbf{Value} \\
\midrule
$N_\textrm{attributes}$ & Number of episode attributes & 2 \\
$|A|, |B| $ & Number of concepts per attribute & [5, 5] \\
$N_\textrm{concept}^A$, $N_\textrm{concept}^B$ & Total size (in neurons) per latent variable & [10, 10] \\
$P(A_i, B_j)$ & Episode probability & $0.5/5$ if $i = j$ ; $0.5/20$ if $i\neq j$ \\
$N_\textrm{swaps}$ & Number of swaps (total noise is $2N_\textrm{swaps}$) & 2 \\
\bottomrule
\end{tabular}
\end{table}
\begin{table}[ht]
\centering
\caption{Summary of default input temporal parameters}
\label{tab:default_time}
\begin{tabular}{lll}
\toprule
\textbf{Variable} & \textbf{Description} & \textbf{Value} \\
\midrule
$T_\textrm{day}$ & Day length & 80 \\
$T_\textrm{episode}$ & Episode length & 5 \\
\bottomrule
\end{tabular}
\end{table}

\begin{table}[ht]
\small
\centering
\caption{Summary of network parameters used in the simulations (default unless specified otherwise).}
\label{tab:default_model}
\begin{tabular}{lp{6cm}l}
\toprule
\textbf{Variable} & \textbf{Description} & \textbf{Value} \\
\midrule
\multicolumn{3}{l}{\textbf{General}} \\
\midrule
$\texttt{duration\_phase\_A}$ & Phase A duration (episodic $\to$ semantic only) & 1000 \\
$\texttt{duration\_phase\_B}$ & Phase B duration (episodic $\to$ semantic + semantic $\to$ episodic) & 1500 \\
$\texttt{sleep\_duration\_episodic}$ & Sleep duration for episodic replay & 10 \\
$\texttt{sleep\_duration\_semantic}$ & Sleep duration for semantic replay & 10 \\
$S_L^\textrm{max}$ & Min semantic load during replay & 1 \\
$S_L^\textrm{min}$ & Max semantic load during input & 2 \\
\midrule
\multicolumn{3}{l}{\textbf{Pattern Completion}} \\
\midrule
$\texttt{mtl\_pattern\_complete\_iterations}$ & Iterations in MTL pattern completion & 10 \\
$\texttt{ctx\_pattern\_complete\_iterations}$ & Iterations in CTX pattern completion & 10 \\
\midrule
\multicolumn{3}{l}{\textbf{Region Parameters}} \\
\midrule
$N_{\text{sen}}$ & Number of SEN neurons & 100 \\
$K^\textrm{SEN}_\textrm{concept}$ & Neurons per concept in SEN & 10 \\
$K^\textrm{SEN}_\textrm{episode}$ & Neurons per episode in SEN & 20 \\
$N^{\text{CTX}}$ & Number of CTX neurons (list denotes sub-regions) & [100, 300] \\
$K^\textrm{CTX}_\textrm{concept}$ & Neurons per concept in CTX (list denotes sub-regions) & [10,  10] \\
$K^\textrm{CTX}_\textrm{episode}$ & Neurons per episode in CTX (list denotes sub-regions) & [20,  10] \\
$b^\text{CTX}$ & CTX immature neurons excitability (list denotes sub-regions) & [0.3, 0.6] \\
$N_{\text{MTL-sensory}}$ & Number of MTL-sensory neurons & 100 \\
$K^\textrm{MTL-sensory}_\textrm{concept}$ & Neurons per concept in MTL-sensory & 10 \\
$K^\textrm{MTL-sensory}_\textrm{episode}$ & Neurons per episode in MTL-sensory & 20 \\
$N_{\text{MTL-semantic}}$ & Number of MTL-semantic neurons & 100 \\
$K^\textrm{MTL-semantic}_\textrm{concept}$ & Neurons per concept in MTL-semantic & 5 \\
$K^\textrm{MTL-semantic}_\textrm{episode}$ & Neurons per episode in MTL-semantic & 10 \\
$b^\text{MTL.semantic}$ & MTL-semantic immature neurons excitability & 0.3 \\
\midrule
\multicolumn{3}{l}{\textbf{Synaptic Parameters}} \\
\midrule
$\texttt{MTL-sensory\_SEN\_projection}$ & MTL-sensory$\leftarrow$SEN random projection & True \\
$W^{\textrm{MTL-sensory}\leftarrow\textrm{SEN}}_\textrm{size}$ & MTL-sensory$\leftarrow$SEN receptive field size & 30 \\
$W^{\text{MTL-sensory}\leftarrow\text{SEN}}_{\text{max-post}}$ & Max incoming MTL-sensory$\leftarrow$SEN weight & 1 \\
$W_\sigma^{\text{CTX}\leftarrow\text{MTL}}$ & CTX$\leftarrow$MTL weight init std &$5\cdot 10^{-4}$ \\
$\lambda^{\text{CTX}\leftarrow\text{MTL}}_{\textrm{slow}}$ & CTX$\leftarrow$MTL learning rate & $5\cdot 10^{-4}$ \\
$W^{\text{CTX}\leftarrow\text{MTL}}_{\text{max-pre}}$ & Max outgoing CTX$\leftarrow$MTL weight & $\infty$ \\
$W^{\text{CTX}\leftarrow\text{MTL}}_{\text{max-post}}$ & Max incoming CTX$\leftarrow$MTL weight & 1 \\
$\lambda^{\text{MTL}\leftarrow\text{MTL}}_0$ & MTL recurrent learning rate & $5\cdot 10^{-3}$ \\
$W^{\text{MTL}\leftarrow\text{MTL}}_{\text{max-pre}}$ & Max outgoing MTL$\leftarrow$MTL weight & $\infty$ \\
$W^{\text{MTL}\leftarrow\text{MTL}}_{\text{max-post}}$ & Max incoming MTL$\leftarrow$MTL weight & $\infty$ \\
$\lambda^{\text{CTX}\leftarrow\text{CTX}}_0$ & CTX recurrent learning rate & $5\cdot 10^{-4}$ \\
$W^{\text{CTX}\leftarrow\text{CTX}}_{\text{max-pre}}$ & Max outgoing CTX$\leftarrow$CTX weight & 1 \\
$W^{\text{CTX}\leftarrow\text{CTX}}_{\text{max-post}}$ & Max incoming CTX$\leftarrow$CTX weight & $\infty$ \\
$W_\sigma^{\text{MTL-semantic}\leftarrow\text{CTX}}$ & MTL-semantic$\leftarrow$CTX weight init std & $1\cdot 10^{-3}$ \\
$\lambda^{\text{MTL-semantic}\leftarrow\text{CTX}}_{\textrm{slow}}$ & MTL-semantic$\leftarrow$CTX learning rate & $5\cdot 10^{-3}$ \\
$W^{\text{MTL-semantic}\leftarrow\text{CTX}}_{\text{max-pre}}$ & Max outgoing MTL-semantic$\leftarrow$CTX weight & $\infty$ \\
$W^{\text{MTL-semantic}\leftarrow\text{CTX}}_{\text{max-post}}$ & Max incoming MTL-semantic$\leftarrow$CTX weight & 1 \\
\bottomrule
\end{tabular}
\end{table}
\begin{table}[ht]
\centering
\caption{Figure 2 parameters (different than or not specified in default)}
\label{tab:figure_2}
\begin{tabular}{lll}
\toprule
\textbf{Variable} & \textbf{Description} & \textbf{Value} \\
\midrule
$\texttt{MTL-sensory\_SEN\_projection}$ & MTL-sensory$\leftarrow$SEN random projection & False \\
$N_\textrm{days}$ & Number of days & 1000 \\
$N_\textrm{swaps}$ & Number of swaps (total noise is $2N_\textrm{swaps}$) & 4 \\
\bottomrule
\end{tabular}
\end{table}

\begin{table}[ht]
\centering
\caption{Figure 3 parameters (different than or not specified in default)}
\label{tab:figure_3}
\begin{tabular}{lll}
\toprule
\textbf{Variable} & \textbf{Description} & \textbf{Value} \\
\midrule
\makecell{$\texttt{MTL-sensory}$\\$\texttt{\_SEN\_projection}$} & MTL-sensory$\leftarrow$SEN random projection & False \\
$N_\textrm{swaps}$ & Number of swaps (total noise is $2N_\textrm{swaps}$) & 4 \\
\midrule
Figs. 3b-3e & & \\
$N_\textrm{days}$ & Number of days & 1500 \\
\midrule
Fig. 3g-3i & & \\
$N_\textrm{Days}$ & Number of days (to get classifier train sample, after model trained) & 100 \\
\midrule
Fig. 3h & & \\
$P(A_i, B_j)$ & Episode probability (train) & 0 if $i=j$ ; 1/20 if $i\neq j$ \\
$P(A_i, B_j)$ & Episode probability (test) & 1/5 if $i=j$ ; 0 else \\
\bottomrule
\end{tabular}
\end{table}


\begin{table}[ht]
\centering
\caption{Figure 4 parameters (different than or not specified in default)}
\label{tab:figure_4}
\begin{tabular}{lll}
\toprule
\textbf{Variable} & \textbf{Description} & \textbf{Value} \\
\midrule
Interleaved (pre-training) \\
$N_\textrm{days}$ & Number of days & 100 \\
$T_\textrm{day}$ & Day length & 200 \\
$T_\textrm{episode}$ & Episode length & 1 \\
\midrule
Blocked (pre-training) \\
$N_\textrm{days}$ & Number of days & 10 (per concept)\\
$T_\textrm{day}$ & Day length & 200  \\
$T_\textrm{episode}$ & Episode length & 1 \\
$P(A_i, B_j)$ & Episode probability & 0.2 if block concept present (0 else) \\
\midrule
Training \\
$N_\textrm{days}$ & Number of days & 1400 \\
\bottomrule
\end{tabular}
\end{table}

\begin{table}[ht]
\centering
\caption{Figure 5 parameters (different than or not specified in default)}
\label{tab:figure_5}
\begin{tabular}{lll}
\toprule
\textbf{Variable} & \textbf{Description} & \textbf{Value} \\
\midrule
Fig. 5c and 5e (Start with best accuracy network from Fig. 4)\\
$N_\textrm{days}$ & Number of days & 1 \\
$T_\textrm{day}$ & Day length & 25 \\
$T_\textrm{episode}$ & Episode length & 5 \\
\midrule
Fig. 5d (Start with best accuracy network from Fig. 4)\\
$N_\textrm{days}$ & Number of days & 1 \\
$T_\textrm{day}$ & Day length & 1 to 5000 (spaced) \\
$T_\textrm{episode}$ & Episode length & 5 \\
\midrule
Fig. 5f \\
\midrule 
Sensory replayed \\
1. Pre-train with blocked (see Table \ref{tab:figure_4})\\
2. 1400 days\\
3. 100 days frozen to test accuracy \\
\midrule
Sensory-semantic replayed \\
$S_L^\textrm{max}$ & \makecell{Max. semantic \\ load during replay} & 1 \\
1. Pre-train with blocked (see Table \ref{tab:figure_4})\\
2. 2000 days \\
3. 100 days frozen to test accuracy \\

\bottomrule
\end{tabular}
\end{table}

\begin{table}[ht]
\centering
\caption{Figure 6 parameters (different than or not specified in default)}
\label{tab:figure_6}
\begin{tabular}{lll}
\toprule
\textbf{Variable} & \textbf{Description} & \textbf{Value} \\
\midrule
Start with best accuracy network from Fig. 5f \\
$N_\textrm{days}$ & Number of days & 500 \\
\bottomrule
\end{tabular}
\end{table}

\subsection*{Reproduction of experimental results}
\subsubsection*{Concept cell formation \shortcite{quiroga_2005, quiroga_2012}}
\paragraph{Experimental goal.}
We quantified concept-cell formation via different measures aimed at quantifying their main defining properties: a high selectivity to individual concepts that generalizes across contexts and metrics. To do so, we measure (i) cross-context linear decodability, (ii) concept selectivity and (iii) normalized mutual information.

\paragraph{Network and latent space.}
We used a pre-trained network (same training regime as Figure 3) and evaluated activity in MTL-sensory and MTL-semantic. As before, episodes were defined as the conjunction of two concepts $(A_i, B_j)$.

\paragraph{Cross-context linear decodability}
For each target concept in $A \cup B$, we built a binary classification task. A training context was sampled by fixing one concept in the opposite attribute (e.g., fixed $B_j$ when decoding $A_i$). Training data then contained exactly two episode types: one positive example (target concept present) and one negative example (target concept absent, replaced by one contrast concept from the same attribute), each presented as one contiguous block. This yields a minimal balanced training set by construction. 
\newline\newline
Testing was performed on randomly sampled episodes under a \emph{context exclusion} rule: all episodes containing the training opposite-context concept were removed. Within this restricted space, the target concept marginal probability was fixed at 0.5 and the remaining probability mass was distributed uniformly across non-target cases. If a sampled test day produced a single class, sampling was repeated until both classes were present.
\newline\newline
For each concept and repeat, we trained two separate decoders (MTL-sensory features and MTL-semantic features) using the same labels. The decoder was a linear logistic classifier (L2-regularized logistic regression). Performance was the held-out test accuracy under context exclusion.

\paragraph{Selectivity analysis on long test streams.}
To characterize concept tuning directly, we ran a long frozen test stream (no consolidation manipulation) with uniform sampling across all 25 latent pairs. We encoded latent identity as one-hot features over all object and scene concepts, and computed neuron-by-latent selectivity as correlation:
\begin{equation}
s_{n\ell}=\mathrm{corr}(x_n, z_\ell),
\end{equation}
where $x_n$ is neuron $n$ activity over samples and $z_\ell$ is latent feature $\ell$. For each neuron, we retained the maximal selectivity $\max_\ell s_{n\ell}$ and compared its distribution across MTL-sensory and MTL-semantic, after thresholding to highly selective units.

\paragraph{Mutual information of most-selective latent.}
For each thresholded neuron, we identified its most-selective latent $\ell^{\*}$ and computed discrete mutual information between binarized neural activity and that latent. To compare across latents with different entropy, we used normalized mutual information:
\begin{equation}
\mathrm{NMI}_n=\frac{I(x_n; z_{\ell^{\*}})}{H(z_{\ell^{\*}})},
\end{equation}
where $I(\cdot;\cdot)$ is mutual information and $H(\cdot)$ is latent entropy. By construction, $\mathrm{NMI}=1$ indicates that neural activity captures all information available about the selected latent.

\begin{table}[ht]
\centering
\caption{Parameters used for concept-cell formation and selectivity analyses.}
\label{tab:concept_cell_parameters}
\begin{tabular}{p{0.58\linewidth}p{0.30\linewidth}}
\toprule
\textbf{Parameter} & \textbf{Value} \\
\midrule
Latent attributes and concepts & 2 attributes ($A$, $B$), 5 concepts each \\
Concept code size (object, scene) & 10 active units, 10 active units \\
Training setup per concept task & 1 fixed opposite context, 1 target concept, 1 contrast concept \\
Training class structure & Two balanced blocks (positive vs negative) \\
Training block duration & 1 timestep per block \\
Training swap noise & 8 swap operations per timestep \\
Test sampling for decoding & 1 day, 200 timesteps, fixed episode duration 5 \\
Test context rule & Exclude all episodes containing training opposite-context concept \\
Target prevalence in test set & 0.5 (remaining mass uniform over non-target cases) \\
Concept repeats & 10 repeats per concept \\
Decoder type & Linear logistic regression (L2 regularization) \\
Regions decoded & MTL-sensory, MTL-semantic \\
Long selectivity test stream & 1 day, 10,000 timesteps, uniform over all latent pairs \\
Selectivity threshold (reported notebook run) & 0.6 \\
Mutual-information readout & $I(x_n;z_{\ell^\*})/H(z_{\ell^\*})$ on thresholded units \\
\bottomrule
\end{tabular}
\end{table}
\subsubsection*{Consolidation promotes the emergence of representational overlap in the hippocampus and medial prefrontal cortex \shortcite{tompary_2017}}

\paragraph{Experimental goal.}
We tested whether offline consolidation increases the similarity of neural representations for episodes that share a common scene context. As in \shortciteA{tompary_2017}, overlap was evaluated before and after sleep-induced consolidation.

\paragraph{Episode set and latent structure.}
Each episode contained one object latent ($A_i$) and one scene latent ($B_j$). We used 128 distinct objects and 4 scenes. Objects were assigned evenly to scenes (32 objects per scene), yielding a deterministic training sequence in which each object was presented with its assigned scene.

\paragraph{Encoding protocol.}
Networks were initialized from scratch and trained for one wake day. Each object--scene episode was presented once, with fixed duration (5 timesteps per episode). Trial-level variability was introduced by swap noise that preserves global sparsity (equal number of active-to-inactive and inactive-to-active flips). This encoding regime produces one noisy realization per episode while keeping category-level structure intact.

\paragraph{Pre- and post-consolidation probing.}
Representational overlap was probed with partial object cues: object information was provided, whereas scene input channels were set to zero. For each cue, MTL cue was pattern completed, fixing the neural activity corresponding to the cue. The completed MTL state was then propagated through the feed-forward/feedback pathway to obtain the final evoked patterns in both MTL and CTX. This probe was performed immediately before sleep (\textit{pre}) and repeated after consolidation (\textit{post}).

\paragraph{Consolidation manipulation.}
Consolidation consisted of one sleep cycle including both replay streams (episodic and semantic replay). Semantic replay was enabled from the beginning of training (i.e., no initial phase with episodic-only replay). Replay durations followed the Figure experiments configuration listed in Table~\ref{tab:overlap_reproduction_parameters}.

\paragraph{Overlap metric and summary statistics.}
For each region (MTL, CTX), we computed a pairwise Pearson correlation matrix across all object-evoked probe patterns. For object $i$, we defined overlap as:
\begin{equation}
\mathrm{OverlapIndex}_i
=
\overline{r}_{i,\mathrm{same\ scene}}
-
\overline{r}_{i,\mathrm{different\ scene}},
\end{equation}
where $\overline{r}_{i,\mathrm{same\ scene}}$ is the mean correlation between object $i$ and other objects assigned to the same scene, and $\overline{r}_{i,\mathrm{different\ scene}}$ is the corresponding mean across objects assigned to different scenes. Positive values indicate stronger scene-dependent clustering. We summarized the distribution of per-object overlap indices by phase and region (mean, standard deviation, and standard error).

\paragraph{Scene-readout control analysis.}
As a complementary readout, we estimated scene identity from the MTL-sensory component of each retrieved pattern by comparing activity across the four scene subspaces and selecting the maximal one. This yielded a per-object scene prediction and an overall decoding accuracy for pre- and post-consolidation phases.

\begin{table}[ht]
\centering
\caption{Parameters used to reproduce consolidation-induced overlap.}
\label{tab:overlap_reproduction_parameters}
\begin{tabular}{p{0.58\linewidth}p{0.30\linewidth}}
\toprule
\textbf{Parameter} & \textbf{Value} \\
\midrule
Number of object concepts ($|A|$) & 128 \\
Number of scene concepts ($|B|$) & 4 \\
Object-to-scene assignment & Balanced (32 objects per scene) \\
Concept code size (object, scene) & 10 active units, 10 active units \\
Sensory/MTL-sensory size & Auto-allocated (1280 object + 40 scene units) \\
Wake training days & 1 \\
Presentations per object per day & 1 \\
Episode duration during encoding & 5 timesteps (fixed) \\
Training noise & 4 swap operations per timestep \\
Cue type for overlap probe & Object-only cue (scene channels silenced) \\
Cue noise at probe & 4 swap operations \\
MTL retrieval during probe & Pattern completion with cue-constrained sensory component \\
Number of consolidation cycles & 1 \\
Replay regime during consolidation & Episodic + semantic replay active \\
Sleep replay duration (episodic stream) & 500 timesteps \\
Sleep replay duration (semantic stream) & 50 timesteps \\
Random seed & 1 \\
\bottomrule
\end{tabular}
\end{table}

\subsubsection*{Neocortical synaptic engrams for remote contextual memories \shortcite{lee_2023}}

\paragraph{Experimental goal and circuit mapping.}
To reproduce the synaptic-engram findings of contextual fear conditioning (\shortciteA{lee_2023}), we mapped cortical and subcortical populations in the model to PFC-like (CTX) and BLA-like (MTL-semantic) representations, respectively. We then quantified how consolidation and extinction reshape synaptic interactions among tagged neuronal subpopulations.

\paragraph{Task design and latent structure.}
Episodes were defined over two latent attributes with 5 concepts each. We interpreted one valence concept ($B_1$) as fear (US) and one context concept ($A_1$) as conditioned context (CS). During conditioning, fear was persistently co-activated with variable context components, establishing a consolidated fear-associated trace. A subsequent recall phase probed reactivation under CS/US-matched input.

\paragraph{Conditioning and recall protocol.}
Conditioning consisted of one wake day followed by sleep consolidation. Recall was then run as a dedicated wake day, again followed by sleep. To match the reference protocol, one sensory channel was randomized within a restricted temporal window during recall, preserving the overall latent identity while introducing within-trial variability. An additional recall pass was used for recall-specific population tagging and downstream synaptic analyses.

\paragraph{Engram and recall-only population definitions.}
Engram populations were defined from mean wake activity during encoding-related windows using the model's sparse activation rule. In CTX, \emph{engram} cells were those active in the conditioning window; in MTL-semantic, engram cells were extracted from the first recall window. We also defined CTX \emph{recall-only} cells as units active during the later recall probe but not active during conditioning. Synaptic-strength distributions were then measured for: (i) CTX engram$\rightarrow$engram vs non-engram$\rightarrow$engram projections, (ii) MTL-semantic engram vs non-engram projections onto CTX engram cells, and (iii) CTX recall-only projections onto CTX engram cells.

\paragraph{Extinction and behavioural proxy.}
Extinction was modeled as repeated CS-context exposure without fear-US latent activation across many days, each day followed by sleep. A fear-expression proxy was computed after each extinction day as the recurrent CTX input onto a pre-identified fear-cell subset, evoked from an early recall-derived cortical state. This yielded a day-by-day trajectory of fear-related drive during extinction.

\paragraph{Post-extinction synaptic reorganization and savings analysis.}
After extinction, we re-estimated CTX engram-related synaptic distributions (engram$\rightarrow$engram and non-engram$\rightarrow$engram) to quantify extinction-associated weakening of engram-specific structure. To evaluate memory savings, we compared recall stability in the extinguished network against a conditioning-only control using pattern completion from matched probe states, with cosine similarity as the stability measure.

\begin{table}[ht]
\centering
\caption{Parameters used for synaptic-engram reproduction.}
\label{tab:synaptic_engrams_parameters}
\begin{tabular}{p{0.58\linewidth}p{0.30\linewidth}}
\toprule
\textbf{Parameter} & \textbf{Value} \\
\midrule
Latent attributes and concepts & 2 attributes, 5 concepts each \\
Concept code size (context, valence) & 10 active units, 10 active units \\
Conditioning schedule & 1 day, 100 timesteps, fixed episode duration 5 \\
Conditioning noise & 4 swap operations per timestep \\
Conditioning latent statistics & Persistent fear-US ($B_1$) with variable context components \\
Recall schedule & 1 day, CS/US-matched probe distribution \\
Recall variability control & Restricted temporal channel randomization (protocol match) \\
Extinction schedule & 1000 days, 20 timesteps/day, fixed episode duration 5 \\
Extinction noise & 4 swap operations per timestep \\
Extinction latent statistics & CS-context presentations without fear-US latent \\
Fear-probe index within recall trajectory & 25\% of recall window \\
Engram tagging basis & Mean wake activity with sparse activation criterion \\
Synaptic comparisons & CTX E$\rightarrow$E vs NE$\rightarrow$E; MTL-semantic E/NE$\rightarrow$CTX E; CTX recall-only$\rightarrow$CTX E \\
Savings metric & Cosine similarity after CTX pattern completion \\
Number of savings seeds & 50 \\
Recording rate for wake indexing & 1 (full-rate recording) \\
Random seed (reported notebook run) & 0 \\
\bottomrule
\end{tabular}
\end{table}

\section*{Acknowledgments}
We would like to thank Rachel Swanson and Colleen J. Gillon, as well as the rest of the members of the Clopath Lab, for discussions and  comments on earlier versions of the manuscript. We gratefully acknowledge the financial support of the Wellcome Trust, the Simons Foundation, the Engineering and Physical Sciences Research Council (EPSRC), and the European Research Council (ERC).
\section*{Declaration of Interests}
The authors declare no competing interests.
\section*{Code Availability}
All code supporting the findings of this study is available at 
\url{https://github.com/albesagonzalez/sensory-semantic-episodes}
\section*{Declaration of generative AI and AI-assisted technologies in the manuscript preparation process}
During the preparation of this work the authors used ChatGPT in order to assist with writing at the sentence level, and coding at the snippet level. After using this tool/service, the authors reviewed and edited the content as needed and take full responsibility for the content of the published article.
\bibliography{Cites.bib}
\clearpage
\section*{Supplementary Material}
% redefine figure numbering for supplements
\renewcommand{\thefigure}{S\arabic{figure}}
\setcounter{figure}{0} % reset counter so it starts at S1

\begin{figure}[h!]
    \centering
    \includegraphics[width=0.65\linewidth]{Figures/Figure_5_supp.png}
    \caption{Same as Fig. \ref{fig:experimental_support_1} but for a network with median accuracy in Fig. \ref{fig:experimental_support_1}f}
    \label{fig:experimental_support_1_supp_1}
\end{figure}


\begin{figure}[h!]
    \centering
    \includegraphics[width=0.95\linewidth]{Figures/Figure_6_supp_1.png}
    \caption{Example recall patterns in Fig. \ref{fig:experimental_support_1}d for only one stored episode (\textbf{a}) and 1000 stored episodes (\textbf{b}). In the limit of infinite stored patterns, a form of over-generalization leads to wrong pattern completion when the presented pattern is of the form $(A_i, B_{j\neq i})$, over-correcting it to the more likely patterns $(A_i, B_i)$ and $(A_j, B_j)$}
    \label{fig:experimental_support_2_supp_1}
\end{figure}

\begin{figure}[h!]
    \centering
    \includegraphics[width=0.95\linewidth]{Figures/Figure_6_supp_2.png}
    \caption{Same as Fig. \ref{fig:experimental_support_2} but for a network with median accuracy in Fig. \ref{fig:experimental_support_1}f}
    \label{fig:experimental_support_2_supp_2}
\end{figure}

\begin{figure}[h!]
    \centering
    \includegraphics[width=0.95\linewidth]{Figures/Figure_7_supp.png}
    \caption{Same as Fig. \ref{fig:higher-order} but starting with a network with median accuracy in Fig. \ref{fig:experimental_support_2}f}
    \label{fig:higher_order_supp}
\end{figure}
\end{document}